% This text has been transcribed in Plain TeX by
% David R. Wilkins
% School of Mathematics, Trinity College, Dublin 2, Ireland
% (dwilkins@maths.tcd.ie)
%
% Trinity College, 2002.

\magnification=\magstep1
\input epsf

\vsize=227 true mm \hsize=170 true mm
   \voffset=-0.4 true mm \hoffset=-5.4 true mm

\def\folio{\ifnum\pageno>0 \number\pageno \else
   \ifnum\pageno<0 \romannumeral-\pageno \else\fi\fi}

\font\Largebf=cmbx10  scaled \magstep2
\font\Largerm=cmr10   scaled \magstep2
\font\largerm=cmr12
\font\largeit=cmti12
\font\largesc=cmcsc10  scaled \magstep1
\font\sc=cmcsc10

\def\overcomma{\mathaccent"7027}
\def\doubleacute{\mathaccent"707D}

\long\def\doublerulebox#1{\vbox{\hrule height0.4pt depth0pt
   \hbox{\vrule width0.4pt\kern1.6pt\vbox{\kern1.6pt\hrule height0.4pt depth0pt
   \hbox{\vrule width0.4pt\kern12pt\vbox{\kern12pt #1\kern12pt}\kern12pt
   \vrule width0.4pt}\hrule height0pt depth0.4pt\kern1.6pt}\kern1.6pt
   \vrule width0.4pt}\hrule height0pt depth0.4pt}}

\pageno=0

\null\vskip72pt

\centerline{\Largebf THE ANALYST}

\vskip24pt

\centerline{\Largebf By}

\vskip24pt

\centerline{\Largebf George Berkeley}

\vfill

\centerline{\largerm Edited by David R. Wilkins}

\vskip 12pt

\centerline{\largerm 2002}

\vskip36pt\eject

\pageno=0

\null

\vfill\eject

\pageno=-3

\null\vskip36pt

\centerline{\Largebf NOTE ON THE TEXT}

\bigskip

This edition is based on the original 1734 first editions
of the {\it Analyst\/} published in London and Dublin,
the copies consulted being those in the Library of
Trinity College, Dublin.  These are the only editions
of the {\it Analyst\/} published in Berkeley's lifetime.
This edition follows closely the 1734 London edition,
some obvious errata being corrected; but a small number
of variant readings from the 1734 Dublin edition have been
introduced.

The 1734 Dublin edition seems to be a reprinting of the
London edition: the punctuation is almost identical, though
there are a number of variations in spelling, and the
capitalization is significantly different.  Moreover there
are a number of errors in this edition, which have
propagated themselves into several later editions of
Berkeley's Collected Works.

The first edition of Berkeley's Collected Works, published
in 1784, reproduced the text of the 1734 Dublin edition,
as did the edition of G.~N.\ Wright, published in 1843.
The errors of the 1734 Dublin edition are all to be found
in these editions, and the edition of G.~N.\ Wright
introduced further errors (e.g., printing $x^n - 1$
and $x^n - 2$ in place of $x^{n-1}$ and $x^{n-2}$ in
Sections XIII.\ and XIV.).

The editions of the Collected Works edited by George Sampson
(London, 1898) and Alexander Campbell Fraser (Oxford, 1901) tend
to follow the 1734 Dublin edition, though correcting some, but
not all, of the more significant errors.  The more modern edition
of Arthur A. Luce (Edinburgh and London, 1951) follows the more
accurate 1734 London edition (though omitting the table of
contents).  The edition of {\it De Motu\/} and the
{\it Analyst\/} edited by Douglas M. Jesseph (Dordrecht, Boston
and London, 1992) is also based on the 1734 London edition, and
follows the punctuation and capitalization of that edition more
closely than other nineteenth and twentieth century editions,
though modernizing the notation in many of the mathematical
formulae.

\bigbreak

This edition reproduces the 18th century mathematical notation
found in the original 1734 editions.  Overlines are used
in place of parentheses for grouping, and powers such as
$x^2$ and $x^3$ are more often written as $xx$ and
$xxx$ respectively, as in the following examples:---

\smallskip

\halign{\qquad $#$\qquad\hfil &&$#$\quad\hfil\cr
\omit\qquad {\it 18th century\/}:\hfil
   &\omit {\it Modern\/}:\hfil \cr\noalign{\vskip3pt}
   \overline{A - {1 \over 2}a} \times \overline{B - {1 \over 2}b} &
   (A - {1 \over 2}a) \times (B - {1 \over 2}b) \cr\noalign{\vskip3pt}
   \mathord{\vbox{\hrule\hbox{$x + o \,$\vrule}}\,}^n &
   (x + o)^n \cr\noalign{\vskip3pt}
   \displaystyle {2yynp + ynnp \over 2yyp} &
   \displaystyle {2y^2 np + yn^2 p \over 2y^2 p} \cr\noalign{\vskip3pt}
   zzz - xxx & z^3 - x^3 \cr}

\vfill\eject

The following errors are common to both the 1734 editions:---

\smallskip

\item{}
{\bf Section~XXIV:}
the letter~`$s$', when used to denote the subsecant $MN$, is
printed sometimes in lowercase and sometimes in uppercase: the
1st, 4th and 5th occurrences in this section being printed in
uppercase.  In this edition, the letter~`$s$' is printed
throughout the section in lowercase, when used to denote the
subsecant.

\smallskip

\item{}
{\bf Section~XLV:}
the sequence of successive fluxions from Isaac Newton's
{\it De Quadratura Curvarum\/} is erroneously given as
`$\doubleacute{z}$.\ $\acute{z}$.\ $z$.\ $\dot{z}$.\
$\ddot{\ddot{z}}$.\ $\dot{\ddot{z}}$', but this has been
corrected in the present edition to read
`$\doubleacute{z}$.\ $\acute{z}$.\ $z$.\ $\dot{z}$.\
$\ddot{z}$.\ $\dot{\ddot{z}}$'.

\bigbreak

In the following instances, this edition follows the
1734 Dublin edition in preference to the 1734 London
edition:---

\smallskip

\item{}
{\bf Sections IX and X:}
in the 1374 London edition, these two paragraphs are run together
to form a single section, labelled IX, and there is no separate
section~X.

\smallskip

\item{}
{\bf Section~XIX:}
the word `Induction' is preceded by `the' in the 1734 London
edition.

\smallskip

\item{}
{\bf Section~L., Query~36:}
the 1734 London edition has `Science of the Principles', where
the 1734 Dublin edition and this edition have `Evidence of the
Principles'.  This variation is discussed in some detail by
Geoffrey Keynes in {\it A bibliography of George Berkeley}
(Oxford, 1976), pp.\ 66--67.  Keynes notes in particular that one
of the copies in the J.~M.\ Keynes Collection, King's College,
Cambridge is inscribed {\it From ye Author\/}, in Berkeley's
hand, on the fly-leaf, and that in this copy the word, printed as
`{\it \ \ ence\/}' has been corrected by hand to `{\it evidence\/}',
leaving no doubt as to Berkeley's intention.  (In the copy of the
1734 London edition in the library of Trinity College, Dublin,
the word `{\it Science\/}' has also been corrected by hand to
`{\it evidence\/}'.)  Keynes also notes in this connection that
the 1734 London edition must have preceded the 1734 Dublin
edition.

\smallskip

\noindent
In addition to the above, there are a number of variant spellings
(listed below) where this edition follows the 1734 Dublin edition
in preference to the 1734 London edition.

\bigbreak

The significant errors in the 1734 Dublin edition (not already
present in the 1734 London edition) are the following:---

\smallskip

\item{}
{\bf Section~X:}
A footnote is given at the beginning of Sect~X.\ that refers
in fact to the final sentence of Sect.~IX.

\smallskip

\item{}
{\bf Section~XVII:}
The date of the letter to Collins is incorrectly given
in the footnote as `Nov.~9, 1679'.  An extract from this
letter, dated Nov.~8, 1676, had been included by William
Jones, in his edition of some of Isaac Newton's mathematical
writings, {\it Analysis per quantitatum series, fluxiones,
ac differentias: cum enumeratione linearum tertii ordinis\/}
(London, 1711).

\smallskip

\item{}
{\bf Section~XXVII:}
the mathematical expression `$o - qo$' is incorrectly given
in the 1734 Dublin edition as `$o - po$'.

\smallskip

\item{}
{\bf Section~XLVI:}
the 1734 Dublin edition omits the phrase `should be it self an
Ordinate'.

\vfill\eject

In addition to the above errors, the following variants
are to be found in the 1734 Dublin edition:---

\smallskip

\item{}
{\bf Table of Contents, IX:}
The 1734 Dublin edition has `Fluxions of a Rectangle' where the 
1734 London edition has `Fluxion of a Rectangle'.

\smallskip

\item{}
{\bf Section~XI:}
the word `very' is omitted before the word `obvious'.

\smallskip

\item{}
{\bf Section~XIV:}
the 1734 Dublin edition has `pre-supposed'
where the 1734 London edition has `presupposeth'.

\smallskip

\item{}
{\bf Section~XVII:}
The 1734 London edition of the {\it Analyst\/}: includes
the following comment, referring to the letter from
Newton to Collins, dated Nov.~8, 1676: `Thus much at least
is plain, that he owned himself satisfied concerning certain
Points, which nevertheless he could not undertake to
demonstrate to others.'  The word `could' is changed to
`would' in the 1734 Dublin edition.  Berkeley's original
wording had been criticized by James Jurin in {\it Geometry
No Friend to Infidelity}, and this may explain the change.
The relevant sentence in Newton's original letter was given
on page~38 of William Jones's 1711 edition of mathematical writings
of Isaac Newton (referred to above), as follows:

\smallskip

\itemitem{}
{\it 
This may seem a bold assertion, because it's hard to say a Figure
may, or may not, be Squar'd, or Compar'd with another; but it's
plain to me by the fountain I draw it from, tho' I will not
undertake to prove it to others.}

\smallskip

\item{}
{\bf Section~XXV:}
The 1734 Dublin edition has `suppposed' where the 1734 London
edition has `supposeth'.

\bigbreak

In the case of the following variant spellings, this edition
follows the 1734 Dublin edition in preference to the 1734
London edition:---

\smallskip

\halign{\qquad #\quad\hfil &&#\quad\hfil\cr
{\it London\/}: &{\it Dublin\/}: &Location in text:\cr
Mathematicalmen &Mathematical-men &Contents, X.\cr
Plains &Planes &Sects.\ III.\ and V.\cr
Polygones &Polygons &Sect.~V.\cr
beside &besides &Sect.~XXX.\cr
compleatly& completely &Sect.~L, Query~53.\cr}

\bigbreak

In the case of the following variant spellings, this edition
follows the 1734 London edition in preference to the 1734
Dublin edition:---

\smallskip

\halign{\qquad #\quad\hfil &&#\quad\hfil\cr
{\it London\/}: &{\it Dublin\/}: &Location in text:\cr
shew'd &shewed &Contents, XIV.\cr
Ancients &Antients &Sect.~III.\cr
Humane &Human &Sect.~IV.\cr
sollicitous &solicitous &Sect~VII.\cr
subtile &subtil &Sects.\ X.\ and XXXII.\cr
Mathematics &Mathematicks &Sect.~XVI., and Sect.~L, Queries 15,
22, 25 and 49.\cr
solved &solv'd &Sect.~XXVI.\cr
Persuasion &Perswasion &Sect.~XLVII.\cr
Public &Publick &Sect.~L.\cr
Trials &Tryals &Sect.~L, Query~34.\cr}

\smallskip

The title page of the 1734 Dublin edition follows that of the
1734 London edition, except that the text at the bottom giving
details of publication has been replaced by the following:---

\vskip6pt

\hbox to \hsize{\hss\vbox{\hsize=396 true pt

\centerline{\largerm DUBLIN:}

\vskip6pt

\noindent
Printed by and for {\sc S.~Fuller} at the {\it Globe\/} in
{\it Meath-street}, and {\it J.~Leathly\/} Bookseller in
{\it Dames-street}.  1734.

\vskip12pt}\hss}

\bigbreak\bigskip

\line{\hfil David R. Wilkins}

\vskip3pt

\line{\hfil Dublin, May 2002}

\vfill\eject

\pageno=0

\null\vfill

\hbox to \hsize{\hss\doublerulebox{\hsize=396 true pt

\vskip12pt

\centerline{\largerm THE}

\vskip6pt

\centerline{\Largerm ANALYST;}

\vskip6pt

\centerline{\largerm OR, A}

\vskip6pt

\centerline{\Largerm DISCOURSE}

\vskip6pt

\centerline{\rm Addressed to an}

\vskip6pt

\centerline{\largerm Infidel \largesc Mathematician.}

\vskip6pt

\centerline{\rm WHEREIN}

\vskip6pt

\noindent
It is examined whether the Object, Principles, and Inferences of
the modern Analysis are more distinctly conceived, or more
evidently deduced, than Religious Mysteries and Points of Faith.

\vskip6pt

\centerline{\vbox{\hrule width\hsize}}

\vskip6pt

\centerline{\rm By the {\sc Author} of
{\it The Minute Philosopher.}}

\vskip6pt

\centerline{\vbox{\hrule width\hsize}}

\vskip6pt

\noindent {\it
First cast out the beam out of thine own Eye; and then shalt thou
see clearly to cast out the mote out of thy brother's eye.}
\hfill S.~Matt.~c.~vii.~v.~5.\quad\null

\vskip6pt

\centerline{\vbox{\hrule width\hsize}}

\vskip6pt
\centerline{\largerm LONDON:}

\vskip6pt

\centerline{\rm Printed for {\sc J.~Tonson} in the {\it Strand}.
   1734.}

\vskip12pt}\hss}

\vfill\eject

\pageno=0

\null

\vfill\eject

\pageno=-9

\null\vskip36pt

\centerline{\largerm THE}

\vskip12pt

\centerline{\Largerm CONTENTS.}

\vskip 12pt

\begingroup\advance\leftskip by 36pt

\item{SECT.~I.}
{\it
Mathematicians presumed to be the great Masters of Reason.  Hence
an undue deference to their decisions where they have no right to
decide.  This one Cause of Infidelity.}

\item{II.}
{\it
Their Principles and Methods to be examined with the same
freedom, which they assume with regard to the Principles and
Mysteries of Religion.  In what Sense and how far Geometry is to
be allowed an Improvement of the Mind.}

\item{III.}
{\it
Fluxions the great Object and Employment of the profound
Geometricians in the present Age.  What these Fluxions are.}

\item{IV.}
{\it
Moments or nascent Increments of flowing Quantities difficult to
conceive.  Fluxions of different Orders.  Second and third
Fluxions obscure Mysteries.}

\item{V.}
{\it
Differences, {\rm i.~e.}\ Increments or Decrements infinitely
small, used by foreign Mathematicians instead of Fluxions or
Velocities of nascent and evanescent Increments.}

\item{VI.}
{\it
Differences of various Orders, {\rm i.~e.}\ Quantities infinitely
less than Quantities infinitely little; and infinitesimal Parts
of infinitesimals of infinitesimals, {\it \&c.}\ without end or
limit.}

\item{VII.}
{\it
Mysteries in faith unjustly objected against by those who admit
them in Science.}

\item{VIII.}
{\it
Modern Analysts supposed by themselves to extend their views even
beyond infinity: Deluded by their own Species or Symbols.}

\item{IX.}
{\it
Method for finding the Fluxion of a Rectangle of two
indeterminate Quantities, shewed to be illegitimate and false.}

\item{X.}
{\it
Implicit Deference of Mathematical-men for the great Author of
Fluxions.  Their earnestness rather to go on fast and far, than
to set out warily and see their way distinctly.}

% 1734 LONDON EDITION: Mathematicalmen
% 1734 DUBLIN EDITION: Mathematical-men

\item{XI.}
{\it
Momentums difficult to comprehend.  No middle Quantity to be
admitted between a finite Quantity and nothing, without admitting
Infinitesimals.}

\item{XII.}
{\it
The Fluxion of any Power of a flowing Quantity.  Lemma premised
in order to examine the method for finding such Fluxion.}

\item{XIII.}
{\it
The rule for the Fluxions of Powers attained by unfair
reasoning.}

\item{XIV.}
{\it
The aforesaid reasoning farther unfolded, and shew'd to be
illogical.}

% 1734 DUBLIN EDITION: shewed

\item{XV.}
{\it
No true Conclusion to be justly drawn by direct consequence from
inconsistent Suppositions.  The same Rules of right reason to be
observed, whether Men argue in Symbols or in Words.}

\item{XVI.}
{\it
An Hypothesis being destroyed, no consequence of such Hypothesis
to be retained.}

\item{XVII.}
{\it
Hard to distinguish between evanescent Increments and
infinitesimal Differences.  Fluxions placed in various Lights.
The great Author, it seems, not satisfied with his own Notions.}

\item{XVIII.}
{\it
Quantities infinitely small supposed and rejected by
{\rm Leibnitz} and his Followers.  No Quantity, according to
them, greater or smaller for the Addition or Subduction of its
Infinitesimal.}

\item{XIX.}
{\it
Conclusions to be proved by the Principles, and not Principles by
the Conclusions.}

\item{XX.}
{\it
The Geometrical Analyst considered as a Logician; and his
Discoveries, not in themselves, but as derived from such
Principles and by such Inferences.}

\item{XXI.}
{\it
A Tangent drawn to the Parabola according to the {\rm calculus
differentialis}.  Truth shewn to be the result of error, and
how.}

\item{XXII.}
{\it
By virtue of a twofold mistake Analysts arrive at Truth, but not
at Science: ignorant how they come at their own Conclusions.}

\item{XXIII.}
{\it
The Conclusion never evident or accurate, in virtue of obscure or
inaccurate Premises.  Finite Quantities might be rejected as well
as Infinitesimals.}

\item{XXIV.}
{\it
The foregoing Doctrine farther illustrated.}

\item{XXV.}
{\it
Sundry Observations thereupon.}

\item{XXVI.}
{\it
Ordinate found from the Area by means of evanescent Increments.}

\item{XXVII.}
{\it
In the foregoing Case the supposed evanescent Increment is really
a finite Quantity, destroyed by an equal Quantity with an
opposite Sign.}

\item{XXVIII.}
{\it
The foregoing Case put generally.  Algebraical Expressions
compared with Geometrical Quantities.}

\item{XXIX.}
{\it
Correspondent Quantities Algebraical and Geometrical equated.
The analysis shewed not to obtain in Infinitesimals, but it must
also obtain in finite Quantities.}

\item{XXX.}
{\it
The getting rid of Quantities by the received Principles, whether
of Fluxions or of Differences, neither good Geometry nor good
Logic.  Fluxions or Velocities, why introduced.}

\item{XXXI.}
{\it
Velocities not to be abstracted from Time and Space: Nor their
Proportions to be investigated or considered exclusively of Time
and Space.}

\item{XXXII.}
{\it
Difficult and obscure Points constitute the Principles of the
modern Analysis, and are the Foundation on which it is built.}

\item{XXXIII.}
{\it
The rational Faculties whether improved by such obscure
Analytics.}

\item{XXXIV.}
{\it
By what inconceivable Steps finite lines are found proportional
to Fluxions.  Mathematical Infidels strain at a Gnat and swallow
a Camel.}

\item{XXXV.}
{\it
Fluxions or Infinitesimals not to be avoided on the received
Principles.  Nice Abstractions and Geometrical Metaphysics.}

\item{XXXVI.}
{\it
Velocities of nascent or evanescent Quantities, whether in
reality understood and signified by finite Lines and Species.}

\item{XXXVII.}
{\it
Signs or Exponents obvious; but Fluxions themselves not so.}

\item{XXXVIII.}
{\it
Fluxions, whether the Velocities with which infinitesimal
Differences are generated?}

\item{XXXIX.}
{\it
Fluxions of Fluxions or second Fluxions, whether to be conceived
as Velocities of Velocities, or rather as Velocities of the
second nascent Increments?}

\item{XL.}
{\it
Fluxions considered, sometimes in one Sense, sometimes in
another: One while in themselves, another in their Exponents:
Hence Confusion and Obscurity.}

\item{XLI.}
{\it
Isochronal Increments, whether finite or nascent, proportional to
their respective Velocities.}

\item{XLII.}
{\it
Time supposed to be divided into Moments: Increments generated in
those Moments: And Velocities proportional to those Increments.}

\item{XLIII.}
{\it
Fluxions, second, third, fourth, {\it \&c.}\ what they are, how
obtained, and how represented.  What Idea of Velocity in a Moment
of Time and Point of Space.}

\item{XLIV.}
{\it
Fluxions of all Orders inconceivable.}

\item{XLV.}
{\it
Signs or Exponents confounded with the Fluxions.}

\item{XLVI.}
{\it
Series of Expressions or of Notes easily contrived.  Whether a
Series, of mere Velocities, or of mere nascent Increments,
corresponding thereunto, be as easily conceived?}

\item{XLVII.}
{\it
Celerities dismissed, and instead thereof Ordinates and Areas
introduced.  Analogies and Expressions useful in the modern
Quadratures, may yet be useless for enabling us to conceive
Fluxions.  No right to apply the Rules without knowledge of the
Principles.}

\item{XLVIII.}
{\it
Metaphysics of modern Analysts most incomprehensible.}

\item{XLIX.}
{\it
Analysts employ'd about notional shadowy Entities.  Their Logics
as exceptionable as their Metaphysics.}

\item{L.}
{\it
Occasion of this Address.  Conclusion.  Queries.}

\par\endgroup

\vfill\eject

\pageno=0

\null

\vfill\eject

\pageno=1

\null\vskip36pt

\centerline{\largerm THE}

\vskip12pt

\centerline{\Largerm ANALYST.}

\vskip12pt

I. \hskip 0pt plus1em
Though I am a Stranger to your Person, yet I am not, Sir, a
Stranger to the Reputation you have acquired, in that branch of
Learning which hath been your peculiar Study; nor to the
Authority that you therefore assume in things foreign to your
Profession, nor to the Abuse that you, and too many more of the
like Character, are known to make of such undue Authority, to the
misleading of unwary Persons in matters of the highest
Concernment, and whereof your mathematical Knowledge can by no
means qualify you to be a competent Judge.  Equity indeed and
good Sense would incline one to disregard the Judgment of Men, in
Points which they have not considered or examined.  But several
who make the loudest Claim to those Qualities, do, nevertheless,
the very thing they would seem to despise, clothing themselves in
the Livery of other Mens Opinions, and putting on a general
deference for the Judgment of you, Gentlemen, who are presumed to
be of all Men the greatest Masters of Reason, to be most
conversant about distinct Ideas, and never to take things on
trust, but always clearly to see your way, as Men whose constant
Employment is the deducing Truth by the justest inference from
the most evident Principles.  With this bias on their Minds, they
submit to your Decisions where you have no right to decide.  And
that this is one short way of making Infidels I am credibly
informed.

\medbreak

II. \hskip 0pt plus1em
Whereas then it is supposed, that you apprehend more distinctly,
consider more closely, infer more justly, conclude more
accurately than other Men, and that you are therefore less
religious because more judicious, I shall claim the privilege of
a Free-Thinker; and take the Liberty to inquire into the Object,
Principles, and Method of Demonstration admitted by the
Mathematicians of the present Age, with the same freedom that you
presume to treat the Principles and Mysteries of Religion; to the
end, that all Men may see what right you have to lead, or what
Encouragement others have to follow you.  It hath been an old
remark that Geometry is an excellent Logic.  And it must be
owned, that when the Definitions are clear; when the Postulata
cannot be refused, nor the Axioms denied; when from the distinct
Contemplation and Comparison of Figures, their Properties are
derived, by a perpetual well-connected chain of Consequences, the
Objects being still kept in view, and the attention ever fixed
upon them; there is acquired a habit of reasoning, close and
exact and methodical: which habit strengthens and sharpens the
Mind, and being transferred to other Subjects, is of general use
in the inquiry after Truth.  But how far this is the case of our
Geometrical Analysts, it may be worth while to consider.

\medbreak

III. \hskip 0pt plus1em
The Method of Fluxions is the general Key, by help whereof the
modern Mathematicians unlock the secrets of Geometry, and
consequently of Nature.  And as it is that which hath enabled
them so remarkably to outgo the Ancients in discovering Theorems
and solving Problems, the exercise and application thereof is
become the main, if not sole, employment of all those who in this
Age pass for profound Geometers.  But whether this Method be
clear or obscure, consistent or repugnant, demonstrative or
precarious, as I shall inquire with the utmost impartiality, so I
submit my inquiry to your own Judgment, and that of every candid
Reader.  Lines are supposed to be generated~\footnote*{Introd. ad
Quadraturam Curvarum.}
by the motion of Points, Planes by the motion of Lines, and
Solids by the motion of Planes.  And whereas Quantities generated
in equal times are greater or lesser, according to the greater or
lesser Velocity, wherewith they increase and are generated, a
Method hath been found to determine Quantities from the
Velocities of their generating Motions.  And such Velocities are
called Fluxions: and the Quantities generated are called flowing
Quantities.  These Fluxions are said to be nearly as the
Increments of the flowing Quantities, generated in the least
equal Particles of time; and to be accurately in the first
Proportion of the nascent, or in the last of the evanescent,
Increments.  Sometimes, instead of Velocities, the momentaneous
Increments or Decrements of undetermined flowing Quantities are
considered, under the Appellation of Moments.

% 1734 DUBLIN EDITION  Planes
% 1734 LONDON EDITION  Plains

\medbreak

IV. \hskip 0pt plus1em
By Moments we are not to understand finite Particles.  These are
said not to be Moments, but Quantities generated from Moments,
which last are only the nascent Principles of finite Quantities.
It is said, that the minutest Errors are not to be neglected in
Mathematics: that the Fluxions are Celerities, not proportional
to the finite Increments though ever so small; but only to the
Moments or nascent Increments, whereof the Proportion alone, and
not the Magnitude, is considered.  And of the aforesaid Fluxions
there be other Fluxions, which Fluxions of Fluxions are called
second Fluxions.  And the Fluxions of these second Fluxions are
called third Fluxions: and so on, fourth, fifth, sixth,
{\it \&c.}\ {\it ad infinitum}.  Now as our Sense is strained and
puzzled with the perception of Objects extremely minute, even so
the Imagination, which Faculty derives from Sense, is very much
strained and puzzled to frame clear Ideas of the least Particles
of time, or the least Increments generated therein: and much more
so to comprehend the Moments, or those Increments of the flowing
Quantities in {\it statu nascenti}, in their very first origin or
beginning to exist, before they become finite Particles.  And it
seems still more difficult, to conceive the abstracted Velocities
of such nascent imperfect Entities.  But the Velocities of the
Velocities, the second, third, fourth, and fifth Velocities,
{\it \&c.}\ exceed, if I mistake not, all Humane Understanding.
The further the Mind analyseth and pursueth these fugitive Ideas,
the more it is lost and bewildered; the Objects, at first
fleeting and minute, soon vanishing out of sight.  Certainly in
any Sense a second or third Fluxion seems an obscure Mystery.
The incipient Celerity of an incipient Celerity, the nascent
Augment of a nascent Augment, {\it i.~e.}\ of a thing which hath
no Magnitude: Take it in which light you please, the clear
Conception of it will, if I mistake not, be found impossible,
whether it be so or no I appeal to the trial of every thinking
Reader.  And if a second Fluxion be inconceivable, what are we to
think of third, fourth, fifth Fluxions, and so onward without
end?

% 1734 LONDON EDITION: Humane Understanding
% 1734 DUBLIN EDITION: Human Understanding

\medbreak

V. \hskip 0pt plus1em
The foreign Mathematicians are supposed by some, even of our own,
to proceed in a manner, less accurate perhaps and geometrical,
yet more intelligible.  Instead of flowing Quantities and their
Fluxions, they consider the variable finite Quantities, as
increasing or diminishing by the continual Addition or Subduction
of infinitely small Quantities.  Instead of the Velocities
wherewith Increments are generated, they consider the Increments
or Decrements themselves, which they call Differences, and which
are supposed to be infinitely small.  The Difference of a Line is
an infinitely little Line; of a Plane an infinitely little Plane.
They suppose finite Quantities to consist of Parts infinitely
little, and Curves to be Polygons, whereof the Sides are
infinitely little, which by the Angles they make one with another
determine the Curvity of the Line.  Now to conceive a Quantity
infinitely small, that is, infinitely less than any sensible or
imaginable Quantity, or any the least finite Magnitude, is, I
confess, above my Capacity.  But to conceive a Part of such
infinitely small Quantity, that shall be still infinitely less
than it, and consequently though multiply'd infinitely shall
never equal the minutest finite Quantity, is, I suspect, an
infinite Difficulty to any Man whatsoever; and will be allowed
such by those who candidly say what they think; provided they
really think and reflect, and do not take things upon trust.

% 1734 DUBLIN EDITION:  Planes
% 1734 LONDON EDITION:  Plains
% 1734 DUBLIN EDITION:  Polygons
% 1734 LONDON EDITION:  Polygones

\medbreak

VI. \hskip 0pt plus1em
And yet in the {\it calculus differentialis}, which Method serves
to all the same Intents and Ends with that of Fluxions, our
modern Analysts are not content to consider only the Differences
of finite Quantities: they also consider the Differences of those
Differences, and the Differences of the Differences of the first
Differences.  And so on {\it ad infinitum}.  That is, they
consider Quantities infinitely less than the least discernible
Quantity; and others infinitely less than those infinitely small
ones; and still others infinitely less than the preceding
Infinitesimals, and so on without end or limit.  Insomuch that we
are to admit an infinite succession of Infinitesimals, each
infinitely less than the foregoing, and infinitely greater than
the following.  As there are first, second, third, fourth, fifth
{\it \&c.}\ Fluxions, so there are Differences, first, second,
third fourth, {\it \&c.}\ in an infinite Progression towards
nothing, which you still approach and never arrive at.  And
(which is most strange) although you should take a Million of
Millions of these Infinitesimals, each whereof is supposed
infinitely greater than some other real Magnitude, and add them
to the least given Quantity, it shall be never the bigger.  For
this is one of the modest {\it postulata\/} of our modern
Mathematicians, and is a Corner-stone or Ground-work of their
Speculations.

\medbreak

VII. \hskip 0pt plus1em
All these Points, I say, are supposed and believed by certain
rigorous Exactors of Evidence in Religion, Men who pretend to
believe no further than they can see.  That Men, who have been
conversant only about clear Points, should with difficulty admit
obscure ones might not seem altogether unaccountable.  But he who
can digest a second or third Fluxion, a second or third
Difference, need not, methinks, be squeamish about any Point in
Divinity.  There is a natural Presumption that Mens Faculties are
made alike.  It is on this Supposition that they attempt to argue
and convince one another.  What, therefore, shall appear
evidently impossible and repugnant to one, may be presumed the
same to another.  But with what appearance of Reason shall any
Man presume to say, that Mysteries may not be Objects of Faith,
at the same time that he himself admits such obscure Mysteries to
be the Object of Science?

\medbreak

VIII. \hskip 0pt plus1em
It must indeed be acknowledged, the modern Mathematicians do not
consider these Points as Mysteries, but as clearly conceived and
mastered by their comprehensive Minds.  They scruple not to say,
that by the help of these new Analytics they can penetrate into
Infinity it self: That they can even extend their Views beyond
Infinity: that their Art comprehends not only Infinite, but
Infinite of Infinite (as they express it) or an Infinity of
Infinites.  But, notwithstanding all these Assertions and
Pretensions, it may be justly questioned whether, as other Men in
other Inquiries are often deceived by Words or Terms, so they
likewise are not wonderfully deceived and deluded by their own
peculiar Signs, Symbols, or Species.  Nothing is easier than to
devise Expressions or Notations for Fluxions and Infinitesimals
of the first, second, third, fourth, and subsequent Orders,
proceeding in the same regular form without end or limit
$\dot{x}$.\ $\ddot{x}$.\ $\dot{\ddot{x}}$.\ $\ddot{\ddot{x}}$.\
{\it \&c.}\ or
$dx$.\ $ddx$.\ $dddx$.\ $ddddx$.\ {\it \&c.}  These Expressions
indeed are clear and distinct, and the Mind finds no difficulty
in conceiving them to be continued beyond any assignable Bounds.
But if we remove the Veil and look underneath, if laying aside
the Expressions we set ourselves attentively to consider the
things themselves, which are supposed to be expressed or marked
thereby, we shall discover much Emptiness, Darkness, and
Confusion; nay, if I mistake not, direct Impossibilities and
Contradictions.  Whether this be the case or no, every thinking
Reader is intreated to examine and judge for himself.

\medbreak

IX. \hskip 0pt plus1em
Having considered the Object, I proceed to consider the
Principles of this new Analysis by Momentums, Fluxions, or
Infinitesimals; wherein if it shall appear that your capital
Points, upon which the rest are supposed to depend, include Error
and false Reasoning; it will then follow that you, who are at a
loss to conduct your selves, cannot with any decency set up for
guides to other Men.  The main Point in the method of Fluxions is
to obtain the Fluxion or Momentum of the Rectangle or Product of
two indeterminate Quantities.  Inasmuch as from thence are
derived Rules for obtaining the Fluxions of all other Products
and Powers; be the Coefficients or the Indexes what they will,
integers or fractions, rational or surd.  Now this fundamental
Point one would think should be very clearly made out,
considering how much is built upon it, and that its Influence
extends throughout the whole Analysis.  But let the Reader judge.
This is given for Demonstration.~\footnote*{Naturalis
Philosophi{\ae} principia mathematica, l.~2.\ lem.~2.}
Suppose the Product or Rectangle $AB$ increased by continual
Motion: and that the momentaneous Increments of the Sides $A$ and
$B$ are $a$ and $b$.  When the Sides $A$ and $B$ were deficient,
or lesser by one half of their Moments, the Rectangle was
$\overline{A - {1 \over 2}a} \times \overline{B - {1 \over 2}b}$,
{\it i.~e.},\
$AB - {1 \over 2}aB - {1 \over 2}bA + {1 \over 4}ab$.
And as soon as the Sides $A$ and $B$ are increased by the other
two halves of their Moments, the Rectangle becomes
$\overline{A + {1 \over 2}a} \times \overline{B + {1 \over 2}b}$
or
$AB + {1 \over 2}aB + {1 \over 2}bA + {1 \over 4}ab$.
From the latter Rectangle subduct the former, and the remaining
Difference will be $aB + bA$.  Therefore the Increment of the
Rectangle generated by the intire Increments $a$ and $b$ is
$aB + bA$.  {\it Q.E.D.}
But it is plain that the direct and true Method to obtain the
Moment or Increment of the Rectangle $AB$, is to take the Sides
as increased by their whole Increments, and so multiply them
together, $A + a$ by $B + b$, the Product whereof
$AB + aB + bA + ab$ is the augmented Rectangle; whence if we
subduct $AB$, the Remainder $aB + bA + ab$ will be the true
Increment of the Rectangle, exceeding that which was obtained by
the former illegitimate and indirect Method by the Quantity $ab$.
And this holds universally be the Quantities $a$ and $b$ what
they will, big or little, Finite or Infinitesimal, Increments,
Moments, or Velocities.  Nor will it avail to say that $ab$ is a
Quantity exceeding small: Since we are told that {\it in rebus
mathematicis errores qu\`{a}m minimi non sunt
contemnendi\/}~\footnote\dag{Introd. ad Quadraturam Curvarum.}.

\medbreak

X. \hskip 0pt plus1em
Such reasoning as this for Demonstration, nothing but the
obscurity of the Subject could have encouraged or induced the
great Author of the Fluxionary Method to put upon his Followers,
and nothing but an implicit deference to Authority could move
them to admit.  The Case indeed is difficult.  There can be
nothing done till you have got rid of the Quantity $ab$.  In
order to this the Notion of Fluxions is shifted: it is placed in
various Lights: Points which should be as clear as first
Principles are puzzled; and Terms which should be steadily used
are ambiguous.  But notwithstanding all this address and skill
the point of getting rid of $ab$ cannot be obtained by legitimate
reasoning.  If a Man by Methods, not geometrical or
demonstrative, shall have satisfied himself of the usefulness of
certain Rules; which he afterwards shall propose to his Disciples
for undoubted Truths; which he undertakes to demonstrate in a
subtile manner, and by the help of nice and intricate Notions; it
is not hard to conceive that such his Disciples may, to save
themselves the trouble of thinking, be inclined to confound the
usefulness of a Rule with the certainty of a Truth, and accept
the one for the other; especially if they are Men accustomed
rather to compute than to think; earnest rather to go on fast and
far, than solicitous to set out warily and see their way
distinctly.

% 1734 LONDON EDITION: Sect 10 appended to Sect 9 without break

\medbreak

XI. \hskip 0pt plus1em
The Points or meer Limits of nascent Lines are undoubtedly equal,
as having no more magnitude one than another, a Limit as such
being no Quantity.  If by a Momentum you mean more than the very
initial Limit, it must be either a finite Quantity or an
Infinitesimal.  But all finite Quantities are expressly excluded
from the Notion of a Momentum.  Therefore the Momentum must be an
Infinitesimal.  And indeed, though much Artifice hath been
employ'd to escape or avoid the admission of Quantities
infinitely small, yet it seems ineffectual.  For ought I see, you
can admit no Quantity as a Medium between a finite Quantity and
nothing, without admitting Infinitesimals.  An Increment
generated in a finite Particle of Time, is it self a finite
Particle; and cannot therefore be a Momentum.  You must therefore
take an Infinitesimal Part of Time wherein to generate your
Momentum.  It is said, the Magnitude of Moments is not
considered:  And yet these same Moments are supposed to be
divided into Parts.  This is not easy to conceive, no more than
it is why we should take Quantities less than $A$ and $B$ in
order to obtain the Increment of $AB$, of which proceeding it
must be owned the final Cause or Motive is very obvious; but it
is not so obvious or easy to explain a just and legitimate Reason
for it, or shew it to be Geometrical.

\medbreak

XII. \hskip 0pt plus1em
From the foregoing Principle so demonstrated, the general Rule
for finding the Fluxion of any Power of a flowing Quantity is
derived~\footnote*{Philosophi{\ae} naturalis principia
Mathematica, lib.~2.\ lem.~2.}.
But, as there seems to have been some inward Scruple or
Consciousness of defect in the foregoing Demonstration, and as
this finding the Fluxion of a given Power is a Point of primary
Importance, it hath therefore been judged proper to demonstrate
the same in a different manner independent of the foregoing
Demonstration.  But whether this other Method be more legitimate
and conclusive than the former, I proceed now to examine; and in
order thereto shall premise the following Lemma.  ``If with a
View to demonstrate any Proposition, a certain Point is supposed,
by virtue of which certain other Points are attained; and such
supposed Point be it self afterwards destroyed or rejected by a
contrary Supposition; in that case, all the other Points,
attained thereby and consequent thereupon, must also be destroyed
and rejected, so as from thence forward to be no more supposed or
applied in the Demonstration.''  This is so plain as to need no
Proof.

\medbreak

XIII. \hskip 0pt plus1em
Now the other Method of obtaining a Rule to find the Fluxion of
any Power is as follows.  Let the Quantity~$x$ flow uniformly,
and be it proposed to find the Fluxion of $x^n$.  In the same
time that $x$ by flowing becomes $x + o$, the Power $x^n$ becomes
$\mathord{\vbox{\hrule\hbox{$x + o \,$\vrule}}\,}^n$,
{\it i.~e.}\ by the Method of infinite Series
$$x^n + nox^{n-1} + {nn-n \over 2} \, oox^{n-2}
   + \hbox{\it \&c.}$$
and the Increments
$$o \hbox{ and }
   nox^{n-1} + {nn-n \over 2} \, oox^{n-2} + \hbox{\it \&c.}$$
are to one another as
$$1 \hbox{ to }
   nx^{n-1} + {nn-n \over 2} \, ox^{n-2} + \hbox{\it \&c.}$$
Let now the Increments vanish, and their last Proportion will be
$1$ to $nx^{n-1}$.  But it should seem that this reasoning is not
fair or conclusive.  For when it is said, let the Increments
vanish, {\it i.~e.}\ let the Increments be nothing, or let there
be no Increments, the former Supposition that the Increments were
something, or that there were Increments, is destroyed, and yet a
Consequence of that Supposition, {\it i.~e.}\ an Expression got
by virtue thereof, is retained.  Which, by the foregoing Lemma,
is a false way of reasoning.  Certainly when we suppose the
Increments to vanish, we must suppose their Proportions, their
Expressions, and every thing else derived from the Supposition of
their Existence to vanish with them.

\medbreak

XIV. \hskip 0pt plus1em
To make this Point plainer, I shall unfold the reasoning, and
propose it in a fuller light to your View.  It amounts therefore
to this, or may in other Words be thus expressed.  I suppose that
the Quantity~$x$ flows, and by flowing is increased, and its
Increment I call $o$, so that by flowing it becomes $x + o$.  And
as $x$ increaseth, it follows that every Power of $x$ is likewise
increased in a due Proportion.  Therefore as $x$ becomes $x + o$,
$x^n$ will become
$\mathord{\vbox{\hrule\hbox{$x + o \,$\vrule}}\,}^n$:
that is, according to the Method of infinite Series,
$$x^n + nox^{n-1} + {nn-n \over 2} \, oox^{n-2}
   + \hbox{\it \&c.}$$
And if from the two augmented Quantities we subduct the Root and
the Power respectively, we shall have remaining the two
Increments, to wit,
$$o \hbox{ and }
   nox^{n-1} + {nn-n \over 2} \, oox^{n-2} + \hbox{\it \&c.}$$
which Increments, being both divided by the common Divisor~$o$,
yield the Quotients
$$1 \hbox{ and }
   nx^{n-1} + {nn-n \over 2} \, ox^{n-2} + \hbox{\it \&c.}$$
which are therefore Exponents of the Ratio of the Increments.
Hitherto I have supposed that $x$ flows, that $x$ hath a real
Increment, that $o$ is something.  And I have proceeded all along
on that Supposition, without which I should not have been able to
have made so much as one single Step.  From that Supposition it
is that I get at the Increment of $x^n$, that I am able to
compare it with the Increment of $x$, and that I find the
Proportion between the two Increments.  I now beg leave to make a
new Supposition contrary to the first, {\it i.~e.}\ I will
suppose that there is no Increment of $x$, or that $o$ is
nothing; which second Supposition destroys my first, and is
inconsistent with it, and therefore with every thing that
supposeth it.  I do nevertheless beg leave to retain $nx^{n-1}$,
which is an Expression obtained in virtue of my first
Supposition, which necessarily presupposeth such Supposition,
and which could not be obtained without it: All which seems a
most inconsistent way of arguing, and such as would not be
allowed of in Divinity.

\medbreak

XV. \hskip 0pt plus1em
Nothing is plainer than that no just Conclusion can be directly
drawn from two inconsistent Suppositions.  You may indeed suppose
any thing possible: But afterwards you may not suppose any thing
that destroys what you first supposed.  Or if you do, you must
begin {\it de novo}.  If therefore you suppose that the Augments
vanish, {\it i.~e.}\ that there are no Augments, you are to begin
again, and see what follows from such Supposition.  But nothing
will follow to your purpose.  You cannot by that means ever
arrive at your Conclusion, or succeed in, what is called by the
celebrated Author, the Investigation of the first or last
Proportions of nascent and evanescent Quantities, by instituting
the Analysis in finite ones.  I repeat it again: You are at
liberty to make any possible Supposition: And you may destroy one
Supposition by another: But then you may not retain the
Consequences, or any part of the Consequences of your first
Supposition so destroyed.  I admit that Signs may be made to
denote either any thing or nothing: And consequently that in the
original Notation $x + o$, $o$ might have signified either an
Increment or nothing.  But then which of these soever you make it
signify, you must argue consistently with such its Signification,
and not proceed upon a double Meaning: which to do were a
manifest Sophism.  Whether you argue in Symbols or in Words, the
Rules of right Reason are still the same.  Nor can it be
supposed, you will plead a Privilege in Mathematics to be exempt
from them.

\medbreak

XVI. \hskip 0pt plus1em
If you assume at first a Quantity increased by nothing, and in
the Expression $x + o$, $o$ stands for nothing, upon this
Supposition as there is no Increment of the Root, so there will
be no Increment of the Power; and consequently there will be none
except the first, of all those Members of the Series constituting
the Power of the Binomial; you will therefore never come at your
Expression of a Fluxion legitimately by such Method.  Hence you
are driven into the fallacious way of proceeding to a certain
Point on the Supposition of an Increment, and then at once
shifting your Supposition to that of no Increment.  There may
seem great Skill in doing this at a certain Point or Period.
Since if this second Supposition had been made before the common
Division by $o$, all had vanished at once, and you must have got
nothing by your Supposition.  Whereas by this Artifice of first
dividing, and then changing your Supposition, you retain $1$ and
$nx^{n-1}$.  But, notwithstanding all this address to cover it,
the fallacy is still the same.  For whether it be done sooner or
later, when once the second Supposition or Assumption is made, in
the same instant the former Assumption and all that you got by it
is destroyed, and goes out together.  And this is universally
true, be the Subject what it will, throughout all the Branches of
humane Knowledge; in any other of which, I believe, Men would
hardly admit such a reasoning as this, which in Mathematics is
accepted for Demonstration.

% 1734 DUBLIN EDITION: Mathematicks

\medbreak

XVII. \hskip 0pt plus1em
It may not be amiss to observe, that the Method for finding the
Fluxion of a Rectangle of two flowing Quantities, as it is set
forth in the Treatise of Quadratures, differs from the
abovementioned taken from the second Book of the Principles, and
is in effect the same with that used in the {\it calculus
differentialis}~\footnote*{Analyse des Infiniment Petits,
part~1.\ prop.~2.}.
For the supposing a Quantity infinitely diminished and therefore
rejecting it, is in effect the rejecting an Infinitesimal; and
indeed it requires a marvellous sharpness of Discernment, to be
able to distinguish between evanescent Increments and
infinitesimal Differences.  It may perhaps be said that the
Quantity being infinitely diminished becomes nothing, and so
nothing is rejected.  But according to the received Principles it
is evident, that no Geometrical Quantity, can by any division or
subdivision whatsoever be exhausted, or reduced to nothing.
Considering the various Arts and Devices used by the great author
of the Fluxionary Method: in how many Lights he placeth his
Fluxions: and in what different ways he attempts to demonstrate
the same Point: one would be inclined to think, he was himself
suspicious of the justness of his own demonstrations; and that he
was not enough pleased with any one notion steadily to adhere to
it.  Thus much at least is plain, that he owned himself satisfied
concerning certain Points, which nevertheless he could not
undertake to demonstrate to others~\footnote*{{\it See
Letter to\/} Collins, Nov.~8, 1676.}.
Whether this satisfaction arose from tentative Methods or
Inductions; which have often been admitted by Mathematicians (for
instance by Dr.~{\it Wallis} in his Arithmetic of Infinites) is
what I shall not pretend to determine.  But, whatever the Case
might have been with respect to the Author, it appears that his
Followers have shewn themselves more eager in applying his
Method, than accurate in examining his Principles.

% 1734 LONDON EDITION: could not undertake to demonstrate to others
% 1734 DUBLIN EDITION: would not undertake to demonstrate to others

% 1734 DUBLIN EDITION: Letter to Collins, Nov. 9, 1679.

\medbreak

XVIII. \hskip 0pt plus1em
It is curious to observe, what subtilty and skill this great
Genius employs to struggle with an insuperable Difficulty; and
through what Labyrinths he endeavours to escape the Doctrine of
Infinitesimals; which as it intrudes upon him whether he will or
no, so it is admitted and embraced by others without the least
repugnance.  {\it Leibnitz\/} and his followers in their {\it
calculus differentialis\/} making no manner of scruple, first to
suppose, and secondly to reject Quantities infinitely small: with
what clearness in the Apprehension and justness in the reasoning,
any thinking Man, who is not prejudiced in favour of those
things, may easily discern.  The Notion or Idea of an
infinitesimal Quantity, as it is an Object simply apprehended by
the Mind, hath been already
considered~\footnote\dag{{\it Sect}.~5 {\it and}~6.}.
I shall now only observe as to the method of getting rid of such
Quantities, that it is done without the least Ceremony.  As in
Fluxions the Point of first importance, and which paves the way
to the rest, is to find the Fluxion of a Product of two
indeterminate Quantities, so in the {\it calculus
differentialis\/} (which Method is supposed to have been borrowed
from the former with some small Alterations) the main Point is to
obtain the difference of such Product.  Now the Rule for this is
got by rejecting the Product or Rectangle of the Differences.
And in general it is supposed, that no Quantity is bigger or
lesser for the Addition or Subduction of its Infinitesimal: and
that consequently no error can arise from such rejection of
Infinitesimals.

\medbreak

XIX. \hskip 0pt plus1em
And yet it should seem that, whatever errors are admitted in the
Premises, proportional errors ought to be apprehended in the
Conclusion, be they finite or infinitesimal: and that therefore
the
$\overcomma{\alpha}\kappa \rho \acute{\iota} \beta \varepsilon
   \iota \alpha$
of Geometry requires nothing should be neglected or rejected.  In
answer to this you will perhaps say, that the Conclusions are
accurately true, and that therefore the Principles and Methods
from whence they are derived must be so too.  But this inverted
way of demonstrating your Principles by your Conclusions, as it
would be peculiar to you Gentlemen, so it is contrary to the
Rules of Logic.  The truth of the Conclusion will not prove
either the Form or the Matter of a Syllogism to be true: inasmuch
as the Illation might have been wrong or the Premises false, and
the Conclusion nevertheless true, though not in virtue of such
Illation or of such Premises.  I say that in every other Science
Men prove their Conclusions by their Principles, and not their
Principles by the Conclusions.  But if in yours you should allow
your selves this unnatural way of proceeding, the Consequence
would be that you must take up with Induction, and bid adieu to
Demonstration.  And if you submit to this, your Authority will no
longer lead the way in Points of Reason and Science.

\medbreak

XX. \hskip 0pt plus1em
I have no Controversy about your Conclusions, but only about your
Logic and Method.  How you demonstrate?  What Objects you are
conversant with, and whether you conceive them clearly?  What
Principles you proceed upon; how sound they may be; and how you
apply them?  It must be remembred that I am not concerned about
the truth of your Theorems, but only about the way of coming at
them; whether it be legitimate or illegitimate, clear or obscure,
scientific or tentative.  To prevent all possibility of your
mistaking me, I beg leave to repeat and insist, that I consider
the Geometrical Analyst as a Logician, {\it i.~e.}\ so far forth
as he reasons and argues; and his Mathematical Conclusions, not
in themselves, but in their Premises; not as true or false,
useful or insignificant, but as derived from such Principles, and
by such Inferences.  And forasmuch as it may perhaps seem an
unaccountable Paradox, that Mathematicians should deduce true
Propositions from false Principles, be right in the Conclusion,
and yet err in the Premises; I shall endeavour particularly to
explain why this may come to pass, and shew how Error may bring
forth Truth, though it cannot bring forth Science.

\medbreak

XXI. \hskip 0pt plus1em
In order therefore to clear up this Point, we will suppose for
instance that a Tangent is to be drawn to a Parabola, and examine
the progress of this Affair, as it is performed by infinitesimal
Differences.

\midinsert
\centerline{\epsfbox{anal_fg.1}}
\endinsert

\noindent
Let $AB$ be a Curve, the Abscisse $AP = x$, the
Ordinate $PB = y$, the Difference of the Abscisse $PM = dx$, the
Difference of the Ordinate $RN = dy$.  Now by supposing the Curve
to be a Polygon, and consequently $BN$, the Increment or
Difference of the Curve, to be a straight Line coincident with
the Tangent, and the differential Triangle $BRN$ to be similar to
the triangle $TPB$ the Subtangent $PT$ is found a fourth
Proportional to $RN : RB : PB$: that is to $dy : dx : y$.  Hence
the Subtangent will be
$\displaystyle {y \,dx \over dy}$.
But herein there is an error arising from the aforementioned
false supposition, whence the value of $PT$ comes out greater
than the Truth: for in reality it is not the Triangle $RNB$ but
$RLB$ which is similar to $PBT$, and therefore (instead of $RN$)
$RL$ should have been the first term of the Proportion,
{\it i.~e.}\ $RN + NL$, {\it i.~e.}\ $dy + z$: whence the true
expression for the Subtangent should have been
$\displaystyle {y\,dx \over dy + z}$.
There was therefore an error of defect in making $dy$ the
divisor: which error was equal to $z$, {\it i.~e.}\ $NL$ the Line
comprehended between the Curve and the Tangent.  Now by the
nature of the Curve $yy=px$, supposing $p$ to be the Parameter,
whence by the rule of Differences $2y\,dy = p\,dx$ and
$\displaystyle dy = {p\,dx \over 2y}$.
But if you multiply $y + dy$ by it self, and retain the whole
Product without rejecting the Square of the Difference, it will
then come out, by substituting the augmented Quantities in the
Equation of the Curve, that
$\displaystyle dy = {p\,dx \over 2y} - {dy\,dy \over 2y}$
truly.  There was therefore an error of excess in making
$\displaystyle dy = {p\,dx \over 2y}$, which followed from the
erroneous Rule of Differences.  And the measure of this second
error is
$\displaystyle {dy\,dy \over 2y} = z$.
Therefore the two errors being equal and contrary destroy each
other; the first error of defect being corrected by a second
error of excess.

\medbreak

XXII. \hskip 0pt plus1em
If you had committed only one error, you would not have come at a
true Solution of the Problem.  But by virtue of a twofold mistake
you arrive, though not at Science, yet at Truth.  For Science it
cannot be called, when you proceed blindfold, and arrive at the
Truth not knowing how or by what means.  To demonstrate that $z$
is equal to
$\displaystyle {dy\,dy \over 2y}$,
let $BR$ or $dx$ be $m$ and $RN$ or $dy$ be $n$.  By the thirty
third Proposition of the first Book of the Conics of {\it
Apollonius}, and from similar Triangles, as $2x$ to $y$ so is $m$
to
$\displaystyle n + z = {my \over 2x}$.
Likewise from the Nature of the Parabola $yy + 2yn + nn = xp +
mp$, and $2yn + nn = mp$: wherefore
$\displaystyle {2yn + nn \over p} = m$:
and because $yy = px$,
$\displaystyle {yy \over p}$
will be equal to $x$.  Therefore substituting these values
instead of $m$ and $x$ we shall have
$$n + z = {my \over 2x} = {2yynp + ynnp \over 2yyp}:$$
{\it i.~e.}
$$n + z = {2yn + nn \over 2y}:$$
which being reduced gives
$$z = {nn \over 2y} = {dy\,dy \over 2y}
   \qquad \hbox{{\it Q.E.D.}}$$

\medbreak

XXIII. \hskip 0pt plus1em
Now I observe in the first place, that the Conclusion comes out
right, not because the rejected Square of $dy$ was infinitely
small; but because this error was compensated by another contrary
and equal error.  I observe in the second place, that whatever is
rejected, be it every so small, if it be real, and consequently
makes a real error in the Premises, it will produce a
proportional real error in the Conclusion.  Your Theorems
therefore cannot be accurately true, nor your Problems accurately
solved, in virtue of Premises, which themselves are not accurate,
it being a rule in Logic that {\it Conclusio sequitur partem
debiliorem}.  Therefore I observe in the third place, that when
the Conclusion is evident and the Premises obscure, or the
Conclusion accurate and the Premises inaccurate, we may safely
pronounce that such Conclusion is neither evident nor accurate,
in virtue of those obscure inaccurate Premises or Principles; but
in virtue of some other Principles which perhaps the Demonstrator
himself never knew or thought of.  I observe in the last place,
that in case the Differences are supposed finite Quantities ever
so great, the Conclusion will nevertheless come out the same:
inasmuch as the rejected Quantities are legitimately thrown out,
not for their smallness, but for another reason, to wit, because
of contrary errors, which destroying each other do upon the whole
cause that nothing is really, though something is apparently
thrown out.  And this Reason holds equally, with respect to
Quantities finite as well as infinitesimal, great as well as
small, a Foot or a Yard long as well as the minutest Increment.

\medbreak

XXIV. \hskip 0pt plus1em
For the fuller illustration of this Point, I shall consider it in
another light, and proceeding in finite Quantities to the
Conclusion, I shall only then make use of one Infinitesimal.

\midinsert
\centerline{\epsfbox{anal_fg.2}}
\endinsert

\noindent
Suppose the straight Line $MQ$ cuts the Curve $AT$ in the points
$R$ and $S$.  Suppose $LR$ a Tangent at the Point~$R$, $AN$ the
Abscisse, $NR$ and $OS$ Ordinates.  Let $AN$ be produced to $O$,
and $RP$ be drawn parallel to $NO$.  Suppose $AN = x$, $NR = y$,
$NO = v$, $PS = z$, the subsecant $MN = s$.  Let the Equation
$y = xx$ express the nature of the Curve: and supposing $y$ and
$x$ increased by their finite Increments, we get
$y + z = xx + 2xv + vv$: whence the former Equation being
subducted there remains $z = 2xv + vv$.  And by reason of similar
Triangles
$$PS : PR :: NR : NM, \quad\hbox{\it i.~e.}\quad z:v::y:s
   = {vy \over z},$$
wherein if for $y$ and $z$ we substitute their values, we get
$${vxx \over 2xv + vv} = s = {xx \over 2x + v}.$$
And supposing $NO$ to be infinitely diminished, the subsecant
$NM$ will in that case coincide with the subtangent $NL$, and $v$
as an Infinitesimal may be rejected, whence it follows that
$$s = NL = {xx \over 2x} = {x \over 2};$$
which is the true value of the Subtangent.  And since this was
obtained by one only error, {\it i.~e.}\ by once rejecting one
only Infinitesimal, it should seem, contrary to what hath been
said, that an infinitesimal Quantity or Difference may be
neglected or thrown away, and the Conclusion nevertheless be
accurately true, although there was no double mistake or
rectifying of one error by another, as in the first Case.  But
if this Point be thoroughly considered, we shall find there is
even here a double mistake, and that one compensates or rectifies
the other.  For in the first place, it was supposed, that when
$NO$ is infinitely diminished or becomes an Infinitesimal, then
the Subsecant $NM$ becomes equal to the Subtangent $NL$.  But
this is a plain mistake, for it is evident, that as a Secant
cannot be a Tangent, so a Subsecant cannot be a Subtangent.  Be
the Difference ever so small, yet still there is a Difference.
And if $NO$ be infinitely small, there will even then be an
infinitely small Difference between $NM$ and $NL$.  Therefore
$NM$ or $s$ was too little for your supposition, (when you
supposed it equal to $NL$) and this error was compensated by a
second error in throwing out $v$, which last error made $s$
bigger than its true value, and in lieu thereof gave the value of
the Subtangent.  This is the true State of the Case, however it
may be disguised.  And to this in reality it amounts, and is at
bottom the same thing, if we should pretend to find the
Subtangent by having first found, from the Equation of the Curve
and similar Triangles, a general Expression for all Subsecants,
and then reducing the Subtangent under this general Rule, by
considering it as the Subsecant when $v$ vanishes or becomes
nothing.

% 1734 EDITIONS: s sometimes erroneously uppercase above

\medbreak

XXV. \hskip 0pt plus1em
Upon the whole I observe, {\it First}, that $v$ can never be
nothing so long as there is a secant.  {\it Secondly}, that the
same Line cannot be both tangent and secant.  {\it Thirdly}, that
when $v$ and $NO$~\footnote*{{\it See the foregoing Figure.}}
vanisheth, $PS$ and $SR$ do also vanish, and with them the
proportionality of the similar Triangles.  Consequently the whole
Expression, which was obtained by means thereof and grounded
thereupon, vanisheth when $v$ vanisheth.  {\it Fourthly}, that
the Method for finding Secants or the Expression of Secants, be
it ever so general, cannot in common sense extend any further
than to all Secants whatsoever: and, as it necessarily supposeth
similar Triangles, it cannot be supposed to take place where
there are not similar Triangles.  {\it Fifthly}, that the
Subsecant will always be less than the Subtangent, and can never
coincide with it; which Coincidence to suppose would be absurd;
for it would be supposing, the same Line at the same time to cut
and not to cut another given Line, which is a manifest
Contradiction, such as subverts the Hypothesis and gives a
Demonstration of its Falshood.  {\it Sixthly}, if this be not
admitted, I demand a Reason why any other apagogical
Demonstration, or Demonstration {\it ad absurdum\/} should be
admitted in Geometry rather than this: Or that some real
Difference be assigned between this and others as such.
{\it Seventhly}, I observe that it is sophistical to suppose $NO$
or $RP$, $PS$, and $SR$ to be finite real Lines in order to form
the Triangle $RPS$, in order to obtain Proportions by similar
Triangles; and afterwards to suppose there are no such Lines, nor
consequently similar Triangles, and nevertheless to retain the
Consequence of the first Supposition, after such Supposition hath
been destroyed by a contrary one.  {\it Eighthly}, That although,
in the present case, by inconsistent Suppositions Truth may be
obtained, yet that such Truth is not demonstrated: That such
Method is not conformable to the Rules of Logic and right Reason:
That, however useful it may be, it must be considered only as a
Presumption, as a Knack, an Art, rather an Artifice, but not a
scientific Demonstration.

\medbreak

XXVI. \hskip 0pt plus1em
The Doctrine premised may be farther illustrated by the following
simple and easy Case, wherein I shall proceed by evanescent
Increments.

\midinsert
\centerline{\epsfbox{anal_fg.3}}
\endinsert

\noindent
Suppose $AB = x$, $BC = y$, $BD = o$, and that $xx$ is
equal to the Area $ABC$: It is proposed to find the Ordinate~$y$
or $BC$.  When $x$ by flowing becomes $x + o$, then $xx$ becomes
$xx + 2xo + oo$: And the Area $ABC$ becomes $ADH$, and the
Increment of $xx$ will be equal to $BDHC$ the Increment of the
Area, {\it i.~e.}\ to $BCFD + CFH$.  And if we suppose the
curvilinear Space $CFH$ to be $qoo$, then $2xo + oo = yo + qoo$,
which divided by $o$ give $2x + o = y + qo$.  And, supposing $o$
to vanish, $2x = y$, in which Case $ACH$ will be a straight Line,
and the Areas $ABC$, $CFH$, Triangles.  Now with regard to this
Reasoning, it hath been already
remarked~\footnote*{{\it Sect}.~12 {\it and\/}~13.\ supra.},
that it is not legitimate or logical to suppose $o$ to vanish,
{\it i.~e.}\ to be nothing, {\it i.~e.}\ that there is no
Increment, unless we reject at the same time with the Increment
it self every Consequence of such Increment,
{\it i.~e.}\ whatsoever could not be obtained but by supposing
such Increment.  It must nevertheless be acknowledged, that the
Problem is rightly solved, and the Conclusion true, to which we
are led by this Method.  It will therefore be asked, how comes it
to pass that the throwing out $o$ is attended with no Error in
the Conclusion?  I answer, the true reason hereof is plainly
this: Because $q$ being Unite, $qo$ is equal to $o$: And
therefore $2x + o - qo = y = 2x$, the equal Quantities $qo$ and
$o$ being destroyed by contrary Signs.

\medbreak

XXVII. \hskip 0pt plus1em
As on the one hand it were absurd to get rid of $o$ by saying,
let me contradict my self: Let me subvert my own Hypothesis: Let
me take it for granted that there is no Increment, at the same
time that I retain a Quantity, which I could never have got at
but by assuming an Increment: So on the other hand it would be
equally wrong to imagine, that in a geometrical Demonstration we
may be allowed to admit any Error, though ever so small, or that
it is possible, in the nature of Things, an accurate Conclusion
should be derived from inaccurate Principles.  Therefore $o$
cannot be thrown out as an Infinitesimal, or upon the Principle
that Infinitesimals may be safely neglected.  But only because it
is destroyed by an equal Quantity with a negative Sign, whence
$o - qo$ is equal to nothing.  And as it is illegitimate to
reduce an Equation, by subducting from one Side a Quantity when
it is not to be destroyed, or when an equal Quantity is not
subducted from the other Side of the Equation: So it must be
allowed a very logical and just Method of arguing, to conclude
that if from Equals either nothing or equal Quantities are
subducted, they shall still remain equal.  And this is a true
Reason why no Error is at last produced by the rejecting of $o$.
Which therefore must not be ascribed to the Doctrine of
Differences, or Infinitesimals, or evanescent Quantities, or
Momentums, or Fluxions.

\medbreak

XXVIII. \hskip 0pt plus1em
Suppose the Case to be general, and that $x^n$ is equal to the
Area $ABC$, whence by the Method of Fluxions the Ordinate is
found $nx^{n-1}$ which we admit for true, and shall inquire how
it is arrived at.  Now if we are content to come at the
Conclusion in a summary way, by supposing that the Ratio of the
Fluxions of $x$ and $x^n$ are found~\footnote*{{\it Sect}.~13.}
to be $1$ and $nx^{n-1}$, and that the Ordinate of the Area is
considered as its Fluxion; we shall not so clearly see our way,
or perceive how the truth comes out, that Method as we have
shewed before being obscure and illogical.  But if we fairly
delineate the Area and its Increment, and divide the latter into
two Parts $BCFD$ and $CFH$~\footnote\dag{{\it See the Figure in
Sect}.~26.},
and proceed regularly by Equations between the algebraical and
geometrical Quantities, the reason of the thing will plainly
appear.  For as $x^n$ is equal to the Area $ABC$, so is the
Increment of $x^n$ equal to the Increment of the Area,
{\it i.~e.}\ to $BDHC$; that is, to say,
$$nox^{n-1} + {nn-n \over 2} \, oox^{n-2} + \hbox{\it \&c.}
   = BDFC + CFH.$$
And only the first Members, on each Side of the Equation being
retained, $nox^{n-1} = BDFC$: and dividing both Sides by $o$ or
$BD$, we shall get $nx^{n-1} = BC$.  Admitting, therefore, that
the curvilinear Space $CFH$ is equal to the rejectaneous Quantity
$${nn-n \over 2} \, oox^{n-2} + \hbox{\it \&c.}$$
and that when this is rejected on one Side, that is rejected on
the other, the Reasoning becomes just and the Conclusion true.
And it is all one whatever Magnitude you allow to $BD$, whether
that of an infinitesimal Difference or a finite Increment ever so
great.  It is therefore plain, that the supposing the
rejectaneous algebraical Quantity to be an infinitely small or
evanescent Quantity, and therefore to be neglected, must have
produced an Error, had it not been for the curvilinear Spaces
being equal thereto, and at the same Time subducted from the
other Part or Side of the Equation agreeably to the Axiom,
{\it If from Equals you subduct Equals, the Remainders will be
equal}.  For those Quantities which by the Analysts are said to
be neglected, or made to vanish, are in reality subducted.  If
therefore the Conclusion be true, it is absolutely necessary that
the finite Space $CFH$ be equal to the Remainder of the Increment
expressed by
$${nn-n \over 2} \, oox^{n-2} \quad \hbox{\it \&c.}$$
equal I say to the finite Remainder of a finite Increment.

% 1734 DUBLIN EDITION: the axiom.  If from equals

\medbreak

XXIX. \hskip 0pt plus1em
Therefore, be the Power what you please, there will arise on one
Side an algebraical Expression, on the other a geometrical
Quantity, each of which naturally divides it self into three
Members: The algebraical or fluxionary Expression, into one
which includes neither the Expression of the Increment of the
Absciss nor of any Power thereof, another which includes the
Expression of the Increment it self, and the third including the
Expression of the Powers of the Increment.  The geometrical
Quantity also or whole increased Area consists of three Parts or
Members, the first of which is the given Area, the second a
Rectangle under the Ordinate and the Increment of the Absciss,
and the third a curvilinear Space.  And, comparing the homologous
or correspondent Members on both Sides, we find that as the first
Member of the Expression is the Expression of the given Area, so
the second Member of the Expression will express the Rectangle or
second Member of the geometrical Quantity; and the third,
containing the Powers of the Increment, will express the
curvilinear Space, or third Member of the geometrical Quantity.
This hint may, perhaps, be further extended and applied to good
purpose, by those who have leisure and curiosity for such
Matters.  The use I make of it is to shew, that the Analysis
cannot obtain in Augments or Differences, but it must also obtain
in finite Quantities, be they ever so great, as was before
observed.

\medbreak

XXX. \hskip 0pt plus1em
It seems therefore upon the whole that we may safely pronounce,
the Conclusion cannot be right, if in order thereto any Quantity
be made to vanish, or be neglected, except that either one Error
is redressed by another; or that secondly, on the same Side of an
Equation equal Quantities are destroyed by contrary Signs, so
that the Quantity we mean to reject is first annihilated; or
lastly, that from opposite Sides equal Quantities are subducted.
And therefore to get rid of Quantities by the received Principles
of Fluxions or of Differences is neither good Geometry nor good
Logic.  When the Augments vanish, the Velocities also vanish.
The Velocities or Fluxions are said to be {\it prim\`{o}\/} and
{\it ultim\`{o}\/}, as the Augments nascent and evanescent.  Take
therefore the {\it Ratio\/} of the evanescent Quantities, it is
the same with that of the Fluxions.  It will therefore answer all
Intents as well.  Why then are Fluxions introduced?  Is it not to
shun or rather to palliate the Use of Quantities infinitely
small?  But we have no Notion whereby to conceive and measure
various Degrees of Velocity, besides Space and Time, or when the
Times are given, besides Space alone.  We have even no Notion of
Velocity prescinded from Time and Space.  When therefore a Point
is supposed to move in given Times, we have no Notion of greater
or lesser Velocities or of Proportions between Velocities, but
only of longer and shorter Lines, and of Proportions between such
Lines generated in equal Parts of Time.

\medbreak

XXXI. \hskip 0pt plus1em
A Point may be the limit of a Line: A Line may be the limit of a
Surface: A Moment may terminate Time.  But how can we conceive a
Velocity by the help of such Limits?  It necessarily implies both
Time and Space, and cannot be conceived without them.  And if the
Velocities of nascent and evanescent Quantities,
{\it i.~e.}\ abstracted from Time and Space, may not be
comprehended, how can we comprehend and demonstrate their
Proportions? Or consider their {\it rationes prim{\ae}\/} and
{\it ultim{\ae}\/}?  For to consider the Proportion or
{\it Ratio\/} of Things implies that such Things have Magnitude:
That such their Magnitudes may be measured, and their Relations
to each other known.  But, as there is no measure of Velocity
except Time and Space, the Proportion of Velocities being only
compounded of the direct Proportion of the Spaces, and the
reciprocal Proportion of the Times; doth it not follow that to
talk of investigating, obtaining, and considering the Proportions
of Velocities, exclusively of Time and Space, is to talk
unintelligibly?

\medbreak

XXXII. \hskip 0pt plus1em
But you will say that, in the use and application of Fluxions,
men do not overstrain their Faculties to a precise Conception of
the abovementioned Velocities, Increments, Infinitesimals, or any
other such like Ideas of a Nature so nice, subtile, and
evanescent.  And therefore you will perhaps maintain, that
Problems may be solved without those inconceivable Suppositions:
and that, consequently, the Doctrine of Fluxions, as to the
practical Part, stands clear of all such Difficulties.  I answer,
that if in the use or application of this Method, those difficult
and obscure Points are not attended to, they are nevertheless
supposed.  They are the Foundations on which the Moderns build,
the Principles on which they proceed, in solving Problems and
discovering Theorems.  It is with the Method of Fluxions as with
all other Methods, which presuppose their respective Principles
and are grounded thereon.  Although the rules may be practised by
Men who neither attend to, nor perhaps know the Principles.  In
like manner, therefore, as a Sailor may practically apply certain
Rules derived from Astronomy and Geometry, the Principles whereof
he doth not understand: And as any ordinary Man may solve divers
numerical Questions, by the vulgar Rules and Operations of
Arithmetic, which he performs and applies without knowing the
Reasons of them: Even so it cannot be denied that you may apply
the Rules of the fluxionary Method: You may compare and reduce
particular Cases to general Forms: You may operate and compute
and solve Problems thereby, not only without an actual Attention
to, or an actual Knowledge of, the Grounds of that Method, and
the Principles whereon it depends, and whence it is deduced, but
even without having ever considered or comprehended them.

% 1734 DUBLIN EDITION: subtil

\medbreak

XXXIII. \hskip 0pt plus1em
But then it must be remembred, that in such Case although you
may pass for an Artist, Computist, or Analyst, yet you may not be
justly esteemed a Man of Science and Demonstration.  Nor should
any Man, in virtue of being conversant in such obscure Analytics,
imagine his rational Faculties to be more improved than those of
other Men, which have been exercised in a different manner, and
on different Subjects; much less erect himself into a Judge and
an Oracle, concerning Matters that have no sort of connexion
with, or dependence on those Species, Symbols or Signs, in the
Management whereof he is so conversant and expert.  As you, who
are a skilful Computist or Analyst, may not therefore be deemed
skilful in Anatomy: or {\it vice versa}, as a Man who can dissect
with Art, may, nevertheless, be ignorant in your Art of
computing: even so you may both, notwithstanding your peculiar
Skill in your respective Arts, be alike unqualified to decide
upon Logic, or Metaphysics, or Ethics, or Religion.  And this
would be true, even admitting that you understood your own
Principles and could demonstrate them.

\medbreak

XXXIV. \hskip 0pt plus1em
If it is said, that Fluxions may be expounded or expressed by
finite Lines proportional to them: Which finite Lines, as they
may be distinctly conceived and known and reasoned upon, so they
may be substituted for the Fluxions, and their mutual Relations
or Proportions be considered as the Proportions of Fluxions: By
which means the Doctrine becomes clear and useful.  I answer that
if, in order to arrive at these finite Lines proportional to the
Fluxions, there be certain Steps made use of which are obscure
and inconceivable, be those finite lines themselves ever so
clearly conceived, it must nevertheless be acknowledged, that
your proceeding is not clear nor your method scientific.
For instance, it is supposed that $AB$ being the Absciss, $BC$
the Ordinate, and $VCH$ a Tangent of the Curve $AC$, $Bb$ or $CE$
the Increment of the Absciss, $Ec$ the Increment of the Ordinate,
which produced meets $VH$ in the Point~$T$, and $Cc$ the
Increment of the Curve.  The right Line $Cc$ being produced to
$K$, there are formed three small Triangles, the Rectilinear
$CEc$, the Mixtilinear $CEc$, and the Rectilinear Triangle $CET$.
It is evident these three Triangles are different from each
other, the Rectilinear $CEc$ being less than the Mixtilinear
$CEc$, whose Sides are the three Increments abovementioned, and
this still less than the Triangle $CET$.  It is supposed that the
Ordinate $bc$ moves into the place $BC$, so that the Point $c$ is
coincident with the Point~$C$; and the right Line $CK$, and
consequently the Curve $Cc$, is coincident with the Tangent $CH$.
In which case the mixtilinear evanescent Triangle $CEc$ will, in
its last form, be similar to the Triangle $CET$: And its
evanescent Sides $CE$, $Ec$ and $Cc$ will be proportional to
$CE$, $ET$ and $CT$ the Sides of the Triangle $CET$.  And
therefore it is concluded, that the Fluxions of the lines $AB$,
$BC$, and $AC$, being in the last Ratio of their evanescent
Increments, are proportional to the Sides of the Triangle $CET$,
or, which is all one, of the Triangle $VBC$ similar
thereunto.~\footnote*{Introd.\ ad Quad.\ Curv.}
It is particularly remarked and insisted on by the great Author,
that the Points $C$ and $c$ must not be distant one from another,
by any the least Interval whatsoever: But that, in order to find
the ultimate Proportions of the Lines $CE$, $Ec$, and $Cc$
({\it i.~e.}\ the Proportions of the Fluxions or Velocities)
expressed by the finite Sides of the Triangle $VBC$, the Points
$C$ and $c$ must be accurately coincident, {\it i.~e.}\ one and
the same.  A Point therefore is considered as a Triangle, or a
Triangle is supposed to be formed in a Point.  Which to conceive
seems quite impossible.  Yet some there are, who, though they
shrink at all other Mysteries, make no difficulty of their own,
who strain at a Gnat and swallow a Camel.

\topinsert
\centerline{\epsfbox{anal_fg.4}}
\endinsert

\medbreak

XXXV. \hskip 0pt plus1em
I know not whether it be worth while to observe, that possibly
some Men may hope to operate by Symbols and Suppositions, in such
sort as to avoid the use of Fluxions, Momentums, and
Infinitesimals after the following manner.  Suppose $x$ to be one
Absciss of a Curve, and $z$ another Absciss of the same Curve.
Suppose also that the respective Areas are $xxx$ and $zzz$: and
that $z - x$ is the Increment of the Absciss, and $zzz - xxx$ the
Increment of the Area, without considering how great, or how
small those Increments may be.  Divide now $zzz - xxx$ by
$z - x$ and the Quotient will be $zz + zx + xx$: and, supposing
that $z$ and $x$ are equal, this same Quotient will be $3xx$
which in that case is the Ordinate, which therefore may be thus
obtained independently of Fluxions and Infinitesimals.  But
herein is a direct Fallacy: for in the first place, it is
supposed that the Abscisses $z$ and $x$ are unequal, without such
supposition no one step could have been made; and in the second
place, it is supposed they are equal; which is a manifest
Inconsistency, and amounts to the same thing that hath been
before considered~\footnote\dag{{\it Sect}.~15.}.
And there is indeed reason to apprehend, that all Attempts for
setting the abstruse and fine Geometry on a right Foundation, and
avoiding the Doctrine of Velocities, Momentums, {\it \&c.}\ will
be found impracticable, till such time as the Object and the End
of Geometry are better understood, than hitherto they seem to
have been.  The great Author of the Method of Fluxions felt this
Difficulty, and therefore he gave in to those nice Abstractions
and Geometrical Metaphysics, without which he saw nothing could
be done on the received Principles; and what in the way of
Demonstration he hath done with them the Reader will judge.  It
must, indeed, be acknowledged, that he used Fluxions, like the
Scaffold of a building, as things to be laid aside or got rid of,
as soon as finite Lines were found proportional to them.  But
then these finite Exponents are found by the help of Fluxions.
Whatever therefore is got by such Exponents and Proportions is to
be ascribed to Fluxions: which must therefore be previously
understood.  And what are these Fluxions?  The Velocities of
evanescent Increments?  And what are these same evanescent
Increments?  They are neither finite Quantities nor Quantities
infinitely small, nor yet nothing.  May we not call them the
Ghosts of departed Quantities?

\medbreak

XXXVI. \hskip 0pt plus1em
Men too often impose on themselves and others, as if they
conceived and understood things expressed by Signs, when in truth
they have no Idea, save only of the very Signs themselves.  And
there are some grounds to apprehend that this may be the present
Case.  The Velocities of evanescent or nascent Quantities are
supposed to be expressed, both by finite Lines of a determinate
Magnitude, and by Algebraical Notes or Signs: but I suspect that
many who, perhaps never having examined the matter, take it for
granted, would upon a narrow scrutiny find it impossible, to
frame any Idea or Notion whatsoever of those Velocities,
exclusive of such finite Quantities and Signs.

\midinsert
\centerline{\epsfbox{anal_fg.5}}
\endinsert

\noindent
Suppose the line
$KP$ described by the Motion of a Point continually accelerated,
and that in equal Particles of time the unequal Parts $KL$, $LM$,
$MN$, $NO$, {\it \&c.}\ are generated.  Suppose also that $a$,
$b$, $c$, $d$, $e$, {\it \&c.}\ denote the Velocities of the
generating Point, at the several Periods of the Parts or
Increments so generated.  It is easy to observe that these
Increments are each proportional to the sum of the Velocities
with which it is described: That, consequently, the several Sums
of the Velocities, generated in equal Parts of Time, may be set
forth by the respective Lines $KL$, $LM$, $MN$,
{\it \&c.}\ generated in the same Times: It is likewise an easy
matter to say, that the last Velocity generated in the first
Particle of Time, may be expressed by the Symbol~$a$, the last in
the second by $b$, the last generated in the third by $c$, and so
on: that $a$ is the Velocity of $LM$ in {\it statu nascenti}, and
$b$, $c$, $d$, $e$, {\it \&c.}\ are the Velocities of the
Increments $MN$, $NO$, $OP$, {\it \&c.}\ in their respective
nascent estates.  You may proceed, and consider these Velocities
themselves as flowing or increasing Quantities, taking the
Velocities of the Velocities, and the Velocities of the
Velocities of the Velocities, {\it i.~e.}\ the first, second,
third {\it \&c.}\ Velocities {\it ad infinitum\/}: which
succeeding Series of Velocities may be thus expressed,
$a$.\ $b - a$.\ $c - 2b + a$.\ $d - 3c + 3b - a$.\ {\it \&c.}\
which you may call by the names of the first, second, third,
fourth Fluxions.  And for an apter Expression you may denote the
variable flowing Line $KL$, $KM$, $KN$, {\it \&c.}\ by the
Letter~$x$; and the first Fluxions by $\dot{x}$, the second by
$\ddot{x}$, the third by $\dot{\ddot{x}}$, and so on {\it ad
infinitum}.

\medbreak

XXXVII. \hskip 0pt plus1em
Nothing is easier than to assign Names, Signs, or Expressions to
these Fluxions, and it is not difficult to compute and operate by
means of such Signs.  But it will be found much more difficult,
to omit the Signs and yet retain in our Minds the things, which
we suppose to be signified by them.  To consider the Exponents,
whether Geometrical, or Algebraical, or Fluxionary, is no
difficult Matter.  But to form a precise Idea of a third Velocity
for instance, in it self and by it self, {\it Hoc opus, hic
labor}.  Nor indeed is it an easy point, to form a clear and
distinct Idea of any Velocity at all, exclusive of and
prescinding from all length of time and space; as also from all
Notes, Signs, or Symbols whatsoever.  This, if I may be allowed
to judge of others by my self, is impossible.  To me it seems
evident, that Measures and Signs are absolutely necessary, in
order to conceive or reason about Velocities; and that,
consequently, when we think to conceive the Velocities, simply
and in themselves, we are deluded by vain Abstractions.

\medbreak

XXXVIII. \hskip 0pt plus1em
It may perhaps be thought by some an easier Method of conceiving
Fluxions, to suppose them the Velocities wherewith the
infinitesimal Differences are generated.  So that the first
Fluxions shall be the Velocities of the first Differences, the
second the Velocities of the second Differences, the third
Fluxions the Velocities of the third Differences, and so on
{\it ad infinitum}.  But not to mention the insurmountable
difficulty of admitting or conceiving Infinitesimals, and
Infinitesimals of Infinitesimals, {\it \&c.}\ it is evident that
this notion of Fluxions would not consist with the great Author's
view; who held that the minutest Quantity ought not to be
neglected, that therefore the Doctrine of Infinitesimal
Differences was not to be admitted in Geometry, and who plainly
appears to have introduced the use of Velocities or Fluxions, on
purpose to exclude or do without them.

\medbreak

XXXIX. \hskip 0pt plus1em
To others it may possibly seem, that we should form a juster Idea
of Fluxions by assuming the finite unequal isochronal Increments
$KL$, $LM$, $MN$, {\it \&c.}\ and considering them in {\it statu
nascenti}, also their Increments in {\it statu nascenti}, and the
nascent Increments of those Increments, and so on, supposing the
first nascent Increments to be proportional to the first Fluxions
or Velocities, the nascent Increments of those Increments to be
proportional to the second Fluxions, the third nascent Increments
to be proportional to the third Fluxions, and so onwards.  And,
as the first Fluxions are the Velocities of the first nascent
Increments, so the second Fluxions may be conceived to be the
Velocities of the second nascent Increments, rather than the
Velocities of Velocities.  But which means the Analogy of
Fluxions may seem better preserved, and the notion rendered more
intelligible.

\medbreak

XL. \hskip 0pt plus1em
And indeed it should seem, that in the way of obtaining the
second or third Fluxion of an Equation, the given Fluxions were
considered rather as Increments than Velocities.  But the
considering them sometimes in one Sense, sometimes in another,
one while in themselves, another in their Exponents, seems to
have occasioned no small share of that Confusion and Obscurity,
which is found in the Doctrine of Fluxions.  It may seem
therefore, that the Notion might be still mended, and that
instead of Fluxions of Fluxions, or Fluxions of Fluxions of
Fluxions, and instead of second, third, or fourth, {\it \&c.}
Fluxions of a given Quantity, it might be more consistent and
less liable to exception to say, the Fluxion of the first nascent
Increment, {\it i.~e.}\ the second Fluxion; the Fluxion of the
second nascent Increment {\it i.~e.}\ the third Fluxion; the
Fluxion of the third nascent Increment, {\it i.~e.}\ the fourth
Fluxion, which Fluxions are conceived respectively proportional,
each to the nascent Principle of the Increment succeeding that
whereof it is the Fluxion.

\medbreak

XLI. \hskip 0pt plus1em
For the more distinct Conception of all which it may be
considered, that if the finite Increment
$LM$~\footnote*{{\it See the foregoing Scheme in Sect}.~36.}
be divided into the Isochronal Parts $Lm$, $mn$, $no$, $oM$; and
the Increment $MN$ divided into the Parts $Mp$, $pq$, $qr$, $rN$
Isochronal to the former; as the whole Increments $LM$, $MN$ are
proportional to the Sums of their describing Velocities, even so
the homologous Particles $Lm$, $Mp$ are also proportional to the
respective accelerated Velocities with which they are described.
And as the Velocity with which $Mp$ is generated, exceeds that
with which $Lm$ was generated, even so the Particle $Mp$ exceeds
the Particle $Lm$.  And in general, as the Isochronal Velocities
describing the Particles of $MN$ exceed the Isochronal Velocities
describing the Particles of $LM$, even so the Particles of the
former exceed the correspondent Particles of the latter.  And
this will hold, be the said Particles ever so small.  $MN$
therefore will exceed $LM$ if they are both taken in their
nascent States: and that excess will be proportional to the
excess of the Velocity~$b$ above the Velocity~$a$.  Hence we may
see that this last account of Fluxions comes, in the upshot, to
the same thing with the first~\footnote\dag{{\it Sect}.~36.}.

\medbreak

XLII. \hskip 0pt plus1em
But notwithstanding what hath been said it must still be
acknowledged, that the finite Particles $Lm$ or $Mp$, though
taken ever so small, are not proportional to the Velocities $a$
and $b$; but each to a Series of Velocities changing every
Moment, or which is the same thing, to an accelerated Velocity,
by which it is generated, during a certain minute Particle of
time: That the nascent beginnings or evanescent endings of finite
Quantities, which are produced in Moments or infinitely small
Parts of Time, are alone proportional to given Velocities: That,
therefore, in order to conceive the first Fluxions, we must
conceive Time divided into Moments, Increments generated in those
Moments, and Velocities proportional to those Increments: That in
order to conceive second and third Fluxions, we must suppose that
the nascent Principles or momentaneous Increments have themselves
also other momentaneous Increments, which are proportional to
their respective generating Velocities: That the Velocities of
these second momentaneous Increments are second Fluxions: those
of their nascent momentaneous Increments third Fluxions.  And so
on {\it ad infinitum}.

\medbreak

XLIII. \hskip 0pt plus1em
By subducting the Increment generated in the first Moment from
that generated in the second, we get the Increment of an
Increment.  And by subducting the Velocity generating in the
first Moment from that generating in the second, we get the
Fluxion of a Fluxion.  In like manner, by subducting the
Difference of the Velocities generating in the two first Moments,
from the excess of the Velocity in the third above that in the
second Moment, we obtain the third Fluxion.  And after the same
Analogy we may proceed to fourth, fifth, sixth Fluxions
{\it \&c.}  And if we call the Velocities of the first, second,
third, fourth Moments, $a$, $b$, $c$, $d$, the Series of Fluxions
will be as above,
$a$.\ $b - a$.\ $c - 2b + a$.\ $d - 3c + 3b - a$.\ {\it ad
infinitum}, {\it i.~e.}\ $\dot{x}$.\ $\ddot{x}$.\
$\dot{\ddot{x}}$.\ $\ddot{\ddot{x}}$.\ {\it ad infinitum}.

\medbreak

XLIV. \hskip 0pt plus1em
Thus Fluxions may be considered in sundry Lights and Shapes,
which seem all equally difficult to conceive.  And indeed, as it
is impossible to conceive Velocity without time or space, without
either finite length or finite
Duration~\footnote\ddag{{\it Sect}.~31.},
it must seem above the powers of Men to comprehend even the first
Fluxions.  And if the first are incomprehensible, what shall we
say of the second and third Fluxions, {\it \&c.}?  He who can
conceive the beginning of a beginning, or the end of an end,
somewhat before the first or after the last, may be perhaps
sharpsighted enough to conceive these things.  But most Men will,
I believe, find it impossible to understand them in any sense
whatever.

\medbreak

XLV. \hskip 0pt plus1em
One would think that Men could not speak too exactly on so nice a
Subject.  And yet, as was before hinted, we may often observe
that the Exponents of Fluxions or Notes representing Fluxions are
confounded with the Fluxions themselves.  Is not this the Case,
when just after the Fluxions of flowing Quantities were said to
be the Celerities of their increasing, and the second Fluxions to
be the mutations of the first Fluxions or Celerities, we are told
that
$\doubleacute{z}$.\ $\acute{z}$.\ $z$.\ $\dot{z}$.\ $\ddot{z}$.\
$\dot{\ddot{z}}$.~\footnote\dag{De Quadratura Curvarum.}
represents a Series of Quantities, whereof each subsequent
Quantity is the Fluxion of the preceding; and each foregoing is a
fluent Quantity having the following one for its Fluxion?

% 1734 EDITIONS: Four dots over penultimate z

\medbreak

XLVI. \hskip 0pt plus1em
Divers Series of Quantities and Expressions, Geometrical and
Algebraical, may be easily conceived, in Lines, in Surfaces, in
Species, to be continued without end or limit.  But it will not
be found so easy to conceive a Series, either of mere Velocities
or of mere nascent Increments, distinct therefrom and
corresponding thereunto.  Some perhaps may be led to think the
Author intended a Series of Ordinates, wherein each Ordinate was
the Fluxion of the preceding and Fluent of the following,
{\it i.~e.}\ that the Fluxion of one Ordinate was it self the
Ordinate of another Curve; and the Fluxion of this last Ordinate
was the Ordinate of yet another Curve; and so on {\it ad
infinitum}.  But who can conceive how the Fluxion (whether
Velocity or nascent Increment) of an Ordinate should be it self
an Ordinate?  Or more than that each preceding Quantity or Fluent
is related to its Subsequent or Fluxion, as the Area of a
curvilinear Figure to its Ordinate; agreeably to what the Author
remarks, that each preceding Quantity in such Series is as the
Area of a curvilinear Figure, whereof the Absciss is $z$, and the
Ordinate is the following Quantity.

% 1734 DUBLIN EDITION: Clause `should be it self an Ordinate' missing

\medbreak

XLVII. \hskip 0pt plus1em
Upon the whole it appears that the Celerities are dismissed, and
instead thereof Areas and Ordinates are introduced.  But however
expedient such Analogies or such Expressions may be found for
facilitating the modern Quadratures, yet we shall not find any
light given us thereby into the original real nature of Fluxions;
or that we are enabled to frame from thence just Ideas of
Fluxions considered in themselves.  In all this the general
ultimate drift of the Author is very clear, but his Principles
are obscure.  But perhaps those Theories of the great Author are
not minutely considered or canvassed by his Disciples; who seem
eager, as was before hinted, rather to operate than to know,
rather to apply his Rules and his Forms, than to understand his
Principles and enter into his Notions.  It is nevertheless
certain, that in order to follow him in his Quadratures, they
must find Fluents from Fluxions; and in order to this, they must
know to find Fluxions from Fluents; and in order to find
Fluxions, they must first know what Fluxions are.  Otherwise they
proceed without Clearness and without Science.  Thus the direct
Method precedes the inverse, and the knowledge of the Principles
is supposed in both.  But as for operating according to Rules,
and by the help of general Forms, whereof the original Principles
and Reasons are not understood, this is to be esteemed merely
technical.  Be the Principles therefore ever so abstruse and
metaphysical, they must be studied by whoever would comprehend
the Doctrine of Fluxions.  Nor can any Geometrician have a right
to apply the Rules of the great Author, without first considering
his metaphysical Notions whence they were derived.  These how
necessary soever in order to Science, which can never be obtained
without a precise, clear, and accurate Conception of the
Principles, are nevertheless by several carelesly passed over;
while the Expressions alone are dwelt on and considered and
treated with great Skill and Management, thence to obtain other
Expressions by Methods, suspicious and indirect (to say the
least) if considered in themselves, however recommended by
Induction and Authority; two Motives which are acknowledged
sufficient to beget a rational Faith and moral Persuasion, but
nothing higher.

\medbreak

XLVIII. \hskip 0pt plus1em
You may possibly hope to evade the Force of all that hath been
said, and to screen false Principles and inconsistent Reasonings,
by a general Pretence that these Objections and Remarks are
Metaphysical.  But this is a vain Pretence.  For the plain Sense
and Truth of what is advanced in the foregoing Remarks, I appeal
to the Understanding of every unprejudiced intelligent Reader.
To the same I appeal, whether the Points remarked upon are not
most incomprehensible Metaphysics.  And Metaphysics not of mine,
but your own.  I would not be understood to infer, that your
Notions are false or vain because they are Metaphysical.  Nothing
is either true or false for that Reason.  Whether a Point be
called Metaphysical or no avails little.  The Question is whether
it be clear or obscure, right or wrong, well or ill-deduced?

\medbreak

XLIX. \hskip 0pt plus1em
Although momentaneous Increments, nascent and evanescent
Quantities, Fluxions and Infinitesimals of all Degrees, are in
truth such shadowy Entities, so difficult to imagine or conceive
distinctly, that (to say the least) they cannot be admitted as
Principles or Objects of clear and accurate Science: and although
this obscurity and incomprehensibility of your Metaphysics had
been alone sufficient, to allay your Pretensions to Evidence; yet
it hath, if I mistake not, been further shewn, that your
Inferences are no more just than your Conceptions are clear, and
that your Logics are as exceptionable as your Metaphysics.  It
should seem therefore upon the whole, that your Conclusions are
not attained by just Reasoning from clear Principles;
consequently that the Employment of modern Analysts, however
useful in mathematical Calculations, and Constructions, doth not
habituate and qualify the Mind to apprehend clearly and infer
justly; and consequently, that you have no right in Virtue of
such Habits, to dictate out of your proper Sphere, beyond which
your Judgment is to pass for no more than that of other Men.

\medbreak

L. \hskip 0pt plus1em
Of a long Time I have suspected, that these modern Analytics were
not scientifical, and gave some Hints thereof to the Public about
twenty five Years ago.  Since which time, I have been diverted by
other Occupations, and imagined I might employ my self better
than in deducing and laying together my Thoughts on so nice a
Subject.  And though of late I have been called upon to make good
my Suggestions; yet, as the Person, who made this Call, doth not
appear to think maturely enough to understand, either those
Metaphysics which he would refute, or Mathematics which he would
patronize, I should have spared my self the trouble of writing
for his Conviction.  Nor should I now have troubled you or my
self with this Address, after so long an Intermission of these
Studies; were it not to prevent, so far as I am able, your
imposing on your self and others in Matters of much higher Moment
and Concern.  And to the end that you may more clearly comprehend
the Force and Design of the foregoing Remarks, and pursue them
still further in your own Meditations, I shall subjoin the
following Queries.

\bigbreak

\begingroup\advance\leftskip by 36pt

\item{\it Query 1.} \hskip 0pt plus1em
Whether the Object of Geometry be not the Proportions of
assignable Extensions?  And whether, there be any need of
considering Quantities either infinitely great or infinitely
small?

\item{\it Qu. 2.} \hskip 0pt plus1em
Whether the end of Geometry be not to measure assignable finite
Extension?  And whether this practical View did not first put Men
on the study of Geometry?

\item{\it Qu. 3.} \hskip 0pt plus1em
Whether the mistaking the Object and End of Geometry hath not
created needless Difficulties, and wrong Pursuits in that
Science?

\item{\it Qu. 4.} \hskip 0pt plus1em
Whether Men may properly be said to proceed in a scientific
Method, without clearly conceiving the Object they are conversant
about, the End proposed, and the Method by which it is pursued?

\item{\it Qu. 5.} \hskip 0pt plus1em
Whether it doth not suffice, that every assignable number of
Parts may be contained in some assignable Magnitude?  And whether
it be not unnecessary, as well as absurd, to suppose that finite
Extension is infinitely divisible?

\item{\it Qu. 6.} \hskip 0pt plus1em
Whether the Diagrams in a Geometrical Demonstration are not to be
considered, as Signs of all possible finite Figures, of all
sensible and imaginable Extensions or Magnitudes of the same
kind?

\item{\it Qu. 7.} \hskip 0pt plus1em
Whether it be possible to free Geometry from insuperable
Difficulties and Absurdities, so long as either the abstract
general Idea of Extension, or absolute external Extension be
supposed its true Object?

\item{\it Qu. 8.} \hskip 0pt plus1em
Whether the Notions of absolute Time, absolute Place, and
absolute Motion be not most abstractedly Metaphysical?  Whether
it be possible for us to measure, compute, or know them?

\item{\it Qu. 9.} \hskip 0pt plus1em
Whether Mathematicians do not engage themselves in Disputes and
Paradoxes, concerning what they neither do nor can conceive?  And
whether the Doctrine of Forces be not a sufficient Proof of
this?~\footnote*{See a {\it Latin\/} treatise, {\it De Motu},
published at {\it London}, in the year 1721.}

\item{\it Qu. 10.} \hskip 0pt plus1em
Whether in Geometry it may not suffice to consider assignable
finite Magnitude, without concerning our selves with Infinity?
And whether it would not be righter to measure large Polygons
having finite Sides, instead of Curves, than to suppose Curves
are Polygons of infinitesimal Sides, a Supposition neither true
nor conceivable?

\item{\it Qu. 11.} \hskip 0pt plus1em
Whether many Points, which are not readily assented to, are not
nevertheless true?  And whether those in the two following
Queries may not be of that Number?

\item{\it Qu. 12.} \hskip 0pt plus1em
Whether it be possible, that we should have had an Idea or Notion
of Extension prior to Motion?  Or whether if a Man had never
perceived Motion, he would ever have known or conceived one thing
to be distant from another?

\item{\it Qu. 13.} \hskip 0pt plus1em
Whether Geometrical Quantity hath coexistent Parts?  And whether
all Quantity be not in a flux as well as Time and Motion?

\item{\it Qu. 14.} \hskip 0pt plus1em
Whether Extension can be supposed an Attribute of a Being
immutable and eternal?

\item{\it Qu. 15.} \hskip 0pt plus1em
Whether to decline examining the Principles, and unravelling the
Methods used in Mathematics, would not shew a bigotry in
Mathematicians?

% 1734 DUBLIN EDITION: Mathematicks

\item{\it Qu. 16.} \hskip 0pt plus1em
Whether certain Maxims do not pass current among Analysts, which
are shocking to good Sense?  And whether the common Assumption
that a finite Quantity divided by nothing is infinite be not of
this Number?

\item{\it Qu. 17.} \hskip 0pt plus1em
Whether the considering Geometrical Diagrams absolutely or in
themselves, rather than as Representatives of all assignable
Magnitudes or Figures of the same kind, be not a principal Cause
of the supposing finite Extension infinitely divisible; and of
all the Difficulties and Absurdities consequent thereupon?

\item{\it Qu. 18.} \hskip 0pt plus1em
Whether from Geometrical Propositions being general, and the
Lines in Diagrams being therefore general Substitutes or
Representatives, it doth not follow that we may not limit or
consider the number of Parts, into which such particular Lines
are divisible?

\item{\it Qu. 19.} \hskip 0pt plus1em
When it is said or implied, that such a certain Line delineated
on Paper contains more than any assignable number of Parts,
whether any more in truth ought to be understood, than that it is
a Sign indifferently representing all finite Lines, be they ever
so great.  In which relative Capacity it contains,
{\it i.~e.}\ stands for more than any assignable number of Parts?
And whether it be not altogether absurd to suppose a finite Line,
considered in it self or in its own positive Nature, should
contain an infinite number of Parts?

\item{\it Qu. 20.} \hskip 0pt plus1em
Whether all Arguments for the infinite Divisibility of finite
Extension do not suppose and imply, either general abstract Ideas
or absolute external Extension to be the Object of Geometry?
And, therefore, whether, along with those Suppositions, such
Arguments also do not cease and vanish?

\item{\it Qu. 21.} \hskip 0pt plus1em
Whether the supposed infinite Divisibility of finite Extension
hath not been a Snare to Mathematicians, and a Thorn in their
Sides?  And whether a Quantity infinitely diminished and a
Quantity infinitely small are not the same thing?

\item{\it Qu. 22.} \hskip 0pt plus1em
Whether it be necessary to consider Velocities of nascent or
evanescent Quantities, or Moments, or Infinitesimals?  And
whether the introducing of Things so inconceivable be not a
reproach to Mathematics?

\item{\it Qu. 23.} \hskip 0pt plus1em
Whether Inconsistencies can be Truths?  Whether Points repugnant
and absurd are to be admitted upon any Subject, or in any
Science?  And whether the use of Infinites ought to be allowed,
as a sufficient Pretext and Apology, for the admitting of such
Points in Geometry?

\item{\it Qu. 24.} \hskip 0pt plus1em
Whether a Quantity be not properly said to be known, when we know
its Proportion to given Quantities?  And whether this Proportion
can be known, but by Expressions or Exponents, either
Geometrical, Algebraical, or Arithmetical?  And whether
Expressions in Lines or Species can be useful but so far forth as
they are reducible to Numbers?

\item{\it Qu. 25.} \hskip 0pt plus1em
Whether the finding out proper Expressions or Notations of
Quantity be not the most general Character and Tendency of the
Mathematics?  And Arithmetical Operation that which limits and
defines their Use?

\item{\it Qu. 26.} \hskip 0pt plus1em
Whether Mathematicians have sufficiently considered the Analogy
and Use of Signs?  And how far the specific limited Nature of
things corresponds thereto?

\item{\it Qu. 27.} \hskip 0pt plus1em
Whether because, in stating a general Case of pure Algebra, we
are at full liberty to make a Character denote, either a positive
or a negative Quantity, or nothing at all, we may therefore in a
geometrical Case, limited by Hypotheses and Reasonings from
particular Properties and Relations of Figures, claim the same
Licence?

\item{\it Qu. 28.} \hskip 0pt plus1em
Whether the Shifting of the Hypothesis, or (as we may call it)
the {\it fallacia Suppositionis\/} be not a Sophism, that far and
wide infects the modern Reasonings, both in the mechanical
Philosophy and in the abstruse and fine Geometry?

\item{\it Qu. 29.} \hskip 0pt plus1em
Whether we can form an Idea or Notion of Velocity distinct from
and exclusive of its Measures, as we can of Heat distinct from
and exclusive of the Degrees on the Thermometer, by which it is
measured?  And whether this be not supposed in the Reasonings of
modern Analysts?

\item{\it Qu. 30.} \hskip 0pt plus1em
Whether Motion can be conceived in a Point of Space?  And if
Motion cannot, whether Velocity can?  And if not, whether a first
or last Velocity can be conceived in a mere Limit, either initial
or final, of the described Space?

\item{\it Qu. 31.} \hskip 0pt plus1em
Where there are no Increments, whether there can be any {\it
Ratio\/} of Increments?  Whether Nothings can be considered as
proportional to real Quantities?  Or whether to talk of their
Proportions be not to talk Nonsense?  Also in what Sense we are
to understand the Proportion of a Surface to a Line, of an Area
to an Ordinate?  And whether Species or Numbers, though properly
expressing Quantities which are not homogeneous, may yet be said
to express their Proportion to each other?

\item{\it Qu. 32.} \hskip 0pt plus1em
Whether if all assignable Circles may be squared, the Circle is
not, to all intents and purposes, squared as well as the
Parabola?  Of whether a parabolical Area can in fact be measured
more accurately than a Circular?

\item{\it Qu. 33.} \hskip 0pt plus1em
Whether it would not be righter to approximate fairly, than to
endeavour at Accuracy by Sophisms?

\item{\it Qu. 34.} \hskip 0pt plus1em
Whether it would not be more decent to proceed by Trials and
Inductions, than to pretend to demonstrate by false Principles?

\item{\it Qu. 35.} \hskip 0pt plus1em
Whether there be not a way of arriving at Truth, although the
Principles are not scientific, nor the Reasoning just?  And
whether such a way ought to be called a Knack or a Science?

\item{\it Qu. 36.} \hskip 0pt plus1em
Whether there can be Science of the Conclusion, where there is
not Evidence of the Principles?  And whether a Man can have
Evidence of the Principles, without understanding them?  And
therefore, whether the Mathematicians of the present Age act like
Men of Science, in taking so much more pains to apply their
Principles, than to understand them?

% 1734 LONDON EDITION: Science of the Principles
% 1734 DUBLIN EDITION: Evidence of the Principles

\item{\it Qu. 37.} \hskip 0pt plus1em
Whether the greatest Genius wrestling with false Principles may
not be foiled?  And whether accurate Quadratures can be obtained
without new {\it Postulata\/} or Assumptions?  And if not,
whether those which are intelligible and consistent ought not to
be preferred to the contrary?
{\it See\/} Sect.\ XXVIII {\it and\/} XXIX.

\item{\it Qu. 38.} \hskip 0pt plus1em
Whether tedious Calculations in Algebra and Fluxions be the
likeliest Method to improve the Mind?  And whether Mens being
accustomed to reason altogether about Mathematical Signs and
Figures, doth not make them at a loss how to reason without them?

\item{\it Qu. 39.} \hskip 0pt plus1em
Whether, whatever readiness Analysts acquire in stating a
Problem, or finding apt Expressions for Mathematical Quantities,
the same doth necessarily infer a proportionable ability in
conceiving and expressing other Matters?

\item{\it Qu. 40.} \hskip 0pt plus1em
Whether it be not a general Case or Rule, that one and the same
Coefficient dividing equal Products gives equal Quotients?  And
yet whether such Coefficient can be interpreted by $o$ or
nothing?  Or whether any one will say, that if the Equation
$2 \times o = 5 \times o$, be divided by $o$, the Quotients on
both Sides are equal?  Whether therefore a Case may not be
general with respect to all Quantities, and yet not extend to
Nothings, or include the Case of Nothing?  And whether the
bringing Nothing under the notion of Quantity may not have
betrayed Men into false Reasoning?

\item{\it Qu. 41.} \hskip 0pt plus1em
Whether in the most general Reasonings about Equalities and
Proportions, Men may not demonstrate as well as in Geometry?
Whether in such Demonstrations, they are not obliged to the same
strict Reasoning as in Geometry?  And whether such their
Reasonings are not deduced from the same Axioms with those in
Geometry?  Whether therefore Algebra be not as truly a Science as
Geometry?

\item{\it Qu. 42.} \hskip 0pt plus1em
Whether Men may not reason in Species as well as in Words?
Whether the same Rules of Logic do not obtain in both Cases?  And
whether we have not a right to expect and demand the same
Evidence in both?

\item{\it Qu. 43.} \hskip 0pt plus1em
Whether an Algebraist, Fluxionist, Geometrician, or Demonstrator
of any kind can expect indulgence for obscure Principles or
incorrect Reasonings?  And whether an Algebraical Note or Species
can at the end of a Process be interpreted in a Sense, which
could not have been substituted for it at the beginning?  Or
whether any particular Supposition can come under a general Case
which doth not consist with the reasoning thereof?

\item{\it Qu. 44.} \hskip 0pt plus1em
Whether the Difference between a mere Computer and a Man of
Science be not, that the one computes on Principles clearly
conceived, and by Rules evidently demonstrated, whereas the other
doth not?

\item{\it Qu. 45.} \hskip 0pt plus1em
Whether, although Geometry be a Science, and Algebra allowed to
be a Science, and the Analytical a most excellent Method, in the
Application nevertheless of the Analysis to Geometry, Men may not
have admitted false Principles and wrong Methods of Reasoning?

\item{\it Qu. 46.} \hskip 0pt plus1em
Whether, although Algebraical Reasonings are admitted to be ever
so just, when confined to Signs or Species as general
Representatives of Quantity, you may not nevertheless fall into
Error, if, when you limit them to stand for particular things,
you do not limit your self to reason consistently with the Nature
of such particular things?  And whether such Error ought to be
imputed to pure Algebra?

\item{\it Qu. 47.} \hskip 0pt plus1em
Whether the View of modern Mathematicians doth not rather seem to
be the coming at an Expression by Artifice, than at the coming at
Science by Demonstration?

\item{\it Qu. 48.} \hskip 0pt plus1em
Whether there may not be sound Metaphysics as well as unsound?
Sound as well as unsound Logic?  And whether the modern Analytics
may not be brought under one of these Denominations, and which?

\item{\it Qu. 49.} \hskip 0pt plus1em
Whether there be not really a {\it Philosophia prima}, a certain
transcendental Science superior to and more extensive than
Mathematics, which it might behove our modern Analysts rather to
learn than despise?

% 1734 DUBLIN EDITION: Mathematicks

\item{\it Qu. 50.} \hskip 0pt plus1em
Whether ever since the recovery of Mathematical Learning, there
have not been perpetual Disputes and Controversies among the
Mathematicians?  And whether this doth not disparage the Evidence
of their Methods?

\item{\it Qu. 51.} \hskip 0pt plus1em
Whether any thing but Metaphysics and Logic can open the Eyes of
Mathematicians and extricate them out of their Difficulties?

\item{\it Qu. 52.} \hskip 0pt plus1em
Whether upon the received Principles a Quantity can by any
Division or Subdivision, though carried ever so far, be reduced
to nothing?

\item{\it Qu. 53.} \hskip 0pt plus1em
Whether if the end of Geometry be Practice, and this Practice be
Measuring, and we measure only assignable Extensions, it will not
follow that unlimited Approximations completely answer the
Intention of Geometry?

% 1734 LONDON EDITION: compleatly

\item{\it Qu. 54.} \hskip 0pt plus1em
Whether the same things which are now done by Infinites may not
be done by finite Quantities?  And whether this would not be a
great Relief to the Imaginations and Understandings of
Mathematical Men?

\item{\it Qu. 55.} \hskip 0pt plus1em
Whether those Philomathematical Physicians, Anatomists, and
Dealers in the Animal Oeconomy, who admit the Doctrine of
Fluxions with an implicit Faith, can with a good grace insult
other Men for believing what they do not comprehend?

\item{\it Qu. 56.} \hskip 0pt plus1em
Whether the Corpuscularian, Experimental, and Mathematical
Philosophy so much cultivated in the last Age, hath not too much
engrossed Mens Attention; some part whereof it might have
usefully employed?

\item{\it Qu. 57.} \hskip 0pt plus1em
Whether from this, and other concurring Causes, the Minds of
speculative Men have not been borne downward, to the debasing and
stupifying of the higher Faculties?  And whether we may not hence
account for that prevailing Narrowness and Bigotry among many who
pass for Men of Science, their Incapacity for things Moral,
Intellectual, or Theological, their Proneness to measure all
Truths by Sense and Experience of animal Life?

\item{\it Qu. 58.} \hskip 0pt plus1em
Whether it be really an Effect of Thinking, that the same Men
admire the great Author for his Fluxions, and deride him for his
Religion?

\item{\it Qu. 59.} \hskip 0pt plus1em
If certain Philosophical Virtuosi of the present Age have no
Religion, whether it can be said to be for want of Faith?

\item{\it Qu. 60.} \hskip 0pt plus1em
Whether it be not a juster way of reasoning, to recommend Points
of Faith from their Effects, than to demonstrate Mathematical
Principles by their Conclusions?

\item{\it Qu. 61.} \hskip 0pt plus1em
Whether it be not less exceptionable to admit Points above Reason
than contrary to Reason?

\item{\it Qu. 62.} \hskip 0pt plus1em
Whether Mysteries may not with better right be allowed of in
Divine Faith, than in Humane Science?

\item{\it Qu. 63.} \hskip 0pt plus1em
Whether such Mathematicians as cry out against Mysteries, have
ever examined their own Principles?

\item{\it Qu. 64.} \hskip 0pt plus1em
Whether Mathematicians, who are so delicate in religious Points,
are strictly scrupulous in their own Science?  Whether they do
not submit to Authority, take things upon Trust, and believe
Points inconceivable?  Whether they have not their Mysteries, and
what is more, their Repugnancies and Contradictions?

\item{\it Qu. 65.} \hskip 0pt plus1em
Whether it might not become Men, who are puzzled and perplexed
about their own Principles, to judge warily, candidly, and
modestly concerning other Matters?

\item{\it Qu. 66.} \hskip 0pt plus1em
Whether the modern Analytics do not furnish a strong {\it
argumentum ad hominem\/} against the Philomathematical Infidels
of these Times?

\item{\it Qu. 67.} \hskip 0pt plus1em
Whether it follows from the abovementioned Remarks, that accurate
and just Reasoning is the peculiar Character of the present Age?
And whether the modern Growth of Infidelity can be ascribed to a
Distinction so truly valuable?

\nobreak\vskip12pt

\centerline{\Largerm FINIS.}

\par\endgroup

\bye
